%&preformat-synopsis
\RequirePackage[l2tabu,orthodox]{nag} % Раскомментировав, можно в логе получать рекомендации относительно правильного
%использования пакетов и предупреждения об устаревших и нерекомендуемых пакетах

% Откомментируйте, чтобы отключить генерацию закладок в pdf
% \PassOptionsToPackage{bookmarks=false}{hyperref}
\documentclass[a5paper,10pt,twoside,openany,article]{memoir} %,draft

%%%%%%%%%%%%%%%%%%%%%%%%%%%%%%%%%%%%%%%%%%%%%%%%%%%%%%%%%%%%%%%%%%%%%%%%%%%%%%%%
%%%% Файл упрощённых настроек шаблона, общих для диссертации и автореферата %%%%
%%%%%%%%%%%%%%%%%%%%%%%%%%%%%%%%%%%%%%%%%%%%%%%%%%%%%%%%%%%%%%%%%%%%%%%%%%%%%%%%

%%% Режим черновика %%%
\makeatletter
\@ifundefined{c@draft}{
  \newcounter{draft}
  \setcounter{draft}{0}  % 0 --- чистовик (максимальное соблюдение ГОСТ)
                         % 1 --- черновик (отклонения от ГОСТ, но быстрая
                         %       сборка итоговых PDF)
}{}
\makeatother

%%% Пометки в тексте %%%
\makeatletter
\@ifundefined{c@showmarkup}{
  \newcounter{showmarkup}
  \setcounter{showmarkup}{1}  % 0 --- скрыть пометки
                              % 1 --- показывать пометки
}{}
\makeatother

%%% Использование в pdflatex шрифтов не по-умолчанию %%%
\makeatletter
\@ifundefined{c@usealtfont}{
  \newcounter{usealtfont}
  \setcounter{usealtfont}{1}    % 0 --- шрифты на базе Computer Modern
                                % 1 --- использовать пакет pscyr, при его
                                %       наличии
                                % 2 --- использовать пакет XCharter, при наличии
                                %       подходящей версии
}{}
\makeatother

%%% Использование в xelatex и lualatex семейств шрифтов %%%
\makeatletter
\@ifundefined{c@fontfamily}{
  \newcounter{fontfamily}
  \setcounter{fontfamily}{2}  % 0 --- CMU семейство. Используется как fallback;
                              % 1 --- Шрифты от MS (Times New Roman и компания)
                              % 2 --- Семейство Liberation
}{}
\makeatother

%%% Библиография %%%
\makeatletter
\@ifundefined{c@bibliosel}{
  \newcounter{bibliosel}
  \setcounter{bibliosel}{1}   % 0 --- встроенная реализация с загрузкой файла
                              %       через движок bibtex8;
                              % 1 --- реализация пакетом biblatex через движок
                              %       biber
}{}
\makeatother

%%% Вывод типов ссылок в библиографии %%%
\makeatletter
\@ifundefined{c@mediadisplay}{
  \newcounter{mediadisplay}
  \setcounter{mediadisplay}{1}   % 0 --- не делать ничего; надписи [Текст] и
                                 %       [Эл. ресурс] будут выводиться только в ссылках с
                                 %       заполненным полем `media`;
                                 % 1 --- автоматически добавлять надпись [Текст] к ссылкам с
                                 %       незаполненным полем `media`; таким образом, у всех
                                 %       источников будет указан тип, что соответствует
                                 %       требованиям ГОСТ
                                 % 2 --- автоматически удалять надписи [Текст], [Эл. Ресурс] и др.;
                                 %       не соответствует ГОСТ
                                 % 3 --- автоматически удалять надпись [Текст];
                                 %       не соответствует ГОСТ
                                 % 4 --- автоматически удалять надпись [Эл. Ресурс];
                                 %       не соответствует ГОСТ
}{}
\makeatother

%%% Предкомпиляция tikz рисунков для ускорения работы %%%
\makeatletter
\@ifundefined{c@imgprecompile}{
  \newcounter{imgprecompile}
  \setcounter{imgprecompile}{0}   % 0 --- без предкомпиляции;
                                  % 1 --- пользоваться предварительно
                                  %       скомпилированными pdf вместо генерации
                                  %       заново из tikz
}{}
\makeatother
          % общие настройки шаблона
%%% Проверка используемого TeX-движка %%%
\newif\ifxetexorluatex   % определяем новый условный оператор (http://tex.stackexchange.com/a/47579)
\ifxetex
    \xetexorluatextrue
\else
    \ifluatex
        \xetexorluatextrue
    \else
        \xetexorluatexfalse
    \fi
\fi

\newif\ifsynopsis           % Условие, проверяющее, что документ --- автореферат

\usepackage{etoolbox}[2015/08/02]   % Для продвинутой проверки разных условий
\providebool{presentation}

\usepackage{comment}    % Позволяет убирать блоки текста (добавляет
                        % окружение comment и команду \excludecomment)

%%% Поля и разметка страницы %%%
\usepackage{pdflscape}  % Для включения альбомных страниц
\usepackage{geometry}   % Для последующего задания полей

%%% Математические пакеты %%%
\usepackage{amsthm,amsmath,amscd}   % Математические дополнения от AMS
\usepackage{amsfonts,amssymb}       % Математические дополнения от AMS
\usepackage{mathtools}              % Добавляет окружение multlined
\usepackage{xfrac}                  % Красивые дроби
\usepackage[
    locale = DE,
    list-separator       = {;\,},
    list-final-separator = {;\,},
    list-pair-separator  = {;\,},
    list-units           = single,
    range-units          = single,
    range-phrase={\text{\ensuremath{-}}},
    % quotient-mode        = fraction, % красивые дроби могут не соответствовать ГОСТ
    fraction-function    = \sfrac,
    separate-uncertainty,
    ]{siunitx}[=v2]                 % Размерности SI
\sisetup{inter-unit-product = \ensuremath{{}\cdot{}}}
% Кириллица в нумерации subequations
% Для правильной работы требуется выполнение сразу после загрузки пакетов
\patchcmd{\subequations}{\def\theequation{\theparentequation\alph{equation}}}
{\def\theequation{\theparentequation\asbuk{equation}}}
{\typeout{subequations patched}}{\typeout{subequations not patched}}

%%%% Установки для размера шрифта 14 pt %%%%
%% Формирование переменных и констант для сравнения (один раз для всех подключаемых файлов)%%
%% должно располагаться до вызова пакета fontspec или polyglossia, потому что они сбивают его работу
\newlength{\curtextsize}
\newlength{\bigtextsize}
\setlength{\bigtextsize}{13.9pt}

\makeatletter
%\show\f@size    % неплохо для отслеживания, но вызывает стопорение процесса,
                 % если документ компилируется без команды  -interaction=nonstopmode
\setlength{\curtextsize}{\f@size pt}
\makeatother

%%% Кодировки и шрифты %%%
\ifxetexorluatex
    \ifpresentation
        \providecommand*\autodot{} % quick fix for polyglossia 1.50
    \fi
    \PassOptionsToPackage{no-math}{fontspec}    % https://tex.stackexchange.com/a/26295/104425
    \usepackage{polyglossia}[2014/05/21]        % Поддержка многоязычности
                                        % (fontspec подгружается автоматически)
\else
   %%% Решение проблемы копирования текста в буфер кракозябрами
    \ifnumequal{\value{usealtfont}}{0}{}{
        \input glyphtounicode.tex
        \input glyphtounicode-cmr.tex %from pdfx package
        \pdfgentounicode=1
    }
    \usepackage{cmap}   % Улучшенный поиск русских слов в полученном pdf-файле
    \ifnumequal{\value{usealtfont}}{2}{}{
        \defaulthyphenchar=127  % Если стоит до fontenc, то переносы
                                % не впишутся в выделяемый текст при
                                % копировании его в буфер обмена
    }
    \usepackage{textcomp}
    \usepackage[T1,T2A]{fontenc}                    % Поддержка русских букв
    \ifnumequal{\value{usealtfont}}{1}{% Используется pscyr, при наличии
        \IfFileExists{pscyr.sty}{\usepackage{pscyr}}{}  % Подключение pscyr
    }{}
    \usepackage[utf8]{inputenc}[2014/04/30]         % Кодировка utf8
    \usepackage[english, russian]{babel}[2014/03/24]% Языки: русский, английский
    \makeatletter\AtBeginDocument{\let\@elt\relax}\makeatother % babel 3.40 fix
    \ifnumequal{\value{usealtfont}}{2}{
        % http://dxdy.ru/post1238763.html#p1238763
        \usepackage[scaled=0.914]{XCharter}[2017/12/19] % Подключение русифицированных шрифтов XCharter
        \usepackage[charter, vvarbb, scaled=1.048]{newtxmath}[2017/12/14]
        \ifpresentation
        \else
            \setDisplayskipStretch{-0.078}
        \fi
    }{}
\fi

%%% Оформление абзацев %%%
\ifpresentation
\else
    \indentafterchapter     % Красная строка после заголовков типа chapter
    \usepackage{indentfirst}
\fi

%%% Цвета %%%
\ifpresentation
\else
    \usepackage[dvipsnames, table, hyperref]{xcolor} % Совместимо с tikz
\fi

%%% Таблицы %%%
\usepackage{longtable,ltcaption} % Длинные таблицы
\usepackage{multirow,makecell}   % Улучшенное форматирование таблиц
\usepackage{tabu, tabulary}      % таблицы с автоматически подбирающейся
                                 % шириной столбцов (tabu обязательно
                                 % до hyperref вызывать)
\makeatletter
%https://github.com/tabu-issues-for-future-maintainer/tabu/issues/26
\@ifpackagelater{longtable}{2020/02/07}{
\def\tabuendlongtrial{%
    \LT@echunk  \global\setbox\LT@gbox \hbox{\unhbox\LT@gbox}\kern\wd\LT@gbox
                \LT@get@widths
}%
}{}
\makeatother

\usepackage{threeparttable}      % автоматический подгон ширины подписи таблицы

%%% Общее форматирование
\usepackage{soulutf8}% Поддержка переносоустойчивых подчёркиваний и зачёркиваний
\usepackage{icomma}  % Запятая в десятичных дробях

%%% Оптимизация расстановки переносов и длины последней строки абзаца
\IfFileExists{impnattypo.sty}{% проверка установленности пакета impnattypo
    \ifluatex
        \ifnumequal{\value{draft}}{1}{% Черновик
            \usepackage[hyphenation, lastparline, nosingleletter, homeoarchy,
            rivers, draft]{impnattypo}
        }{% Чистовик
            \usepackage[hyphenation, lastparline, nosingleletter]{impnattypo}
        }
    \else
        \usepackage[hyphenation, lastparline]{impnattypo}
    \fi
}{}

%% Векторная графика

\usepackage{tikz}                   % Продвинутый пакет векторной графики
\usetikzlibrary{chains}             % Для примера tikz рисунка
\usetikzlibrary{shapes.geometric}   % Для примера tikz рисунка
\usetikzlibrary{shapes.symbols}     % Для примера tikz рисунка
\usetikzlibrary{arrows}             % Для примера tikz рисунка

\usepackage[european,cuteinductors]{circuitikz} % Электрические схемы
\usepackage{pgfplots}                           % Графики
\pgfplotsset{compat=newest}
\usepgfplotslibrary{groupplots,units}
\pgfkeys{/pgf/number format/.cd,use comma,1000 sep={}} % форматирование чисел в графиках

%%% Гиперссылки %%%
\ifxetexorluatex
    \let\CYRDZE\relax
\fi
\usepackage{hyperref}[2012/11/06]

%%% Изображения %%%
\usepackage{graphicx}[2014/04/25]   % Подключаем пакет работы с графикой
\usepackage{caption}                % Подписи рисунков и таблиц
\usepackage{subcaption}             % Подписи подрисунков и подтаблиц
\usepackage{pdfpages}               % Добавление внешних pdf файлов

%%% Счётчики %%%
\usepackage{aliascnt}
\usepackage[figure,table]{totalcount}   % Счётчик рисунков и таблиц
\usepackage{totcount}   % Пакет создания счётчиков на основе последнего номера
                        % подсчитываемого элемента (может требовать дважды
                        % компилировать документ)
\usepackage{totpages}   % Счётчик страниц, совместимый с hyperref (ссылается
                        % на номер последней страницы). Желательно ставить
                        % последним пакетом в преамбуле

%%% Продвинутое управление групповыми ссылками (пока только формулами) %%%
\ifpresentation
\else
    \usepackage[russian]{cleveref} % cleveref имеет сложности со считыванием
    % языка из babel. Такое решение русификации вывода выбрано вместо
    % определения в documentclass из опасности что-то лишнее передать во все
    % остальные пакеты, включая библиографию.

    % Добавление возможности использования пробелов в \labelcref
    % https://tex.stackexchange.com/a/340502/104425
    \usepackage{kvsetkeys}
    \makeatletter
    \let\org@@cref\@cref
    \renewcommand*{\@cref}[2]{%
        \edef\process@me{%
            \noexpand\org@@cref{#1}{\zap@space#2 \@empty}%
        }\process@me
    }
    \makeatother
\fi

\usepackage{placeins} % для \FloatBarrier

%%% Цитата, не приводимая в автореферате:
% возможно, актуальна только для biblatex
%\newcommand{\citeinsynopsis}[1]{\ifsynopsis\else ~\cite{#1} \fi}

% если текущий процесс запущен библиотекой tikz-external, то прекомпиляция должна быть включена
\ifdefined\tikzexternalrealjob
    \setcounter{imgprecompile}{1}
\fi

\ifnumequal{\value{imgprecompile}}{1}{% Только если у нас включена предкомпиляция
    \usetikzlibrary{external}   % подключение возможности предкомпиляции
    \tikzexternalize[prefix=images/cache/,optimize command away=\includepdf] % activate! % здесь можно указать отдельную папку для скомпилированных файлов
    \ifxetex
        \tikzset{external/up to date check={diff}}
    \fi
}{}
       % Пакеты общие для диссертации и автореферата
\synopsistrue                 % Этот документ --- автореферат
\input{Synopsis/synpackages}  % Пакеты для автореферата
\input{Synopsis/userpackages} % Пакеты для специфических пользовательских задач

% Новые переменные, которые могут использоваться во всём проекте
% ГОСТ 7.0.11-2011
% 9.2 Оформление текста автореферата диссертации
% 9.2.1 Общая характеристика работы включает в себя следующие основные структурные
% элементы:
% актуальность темы исследования;
\newcommand{\actualityTXT}{Актуальность темы.}
% степень ее разработанности;
\newcommand{\progressTXT}{Степень разработанности темы.}
% цели и задачи;
\newcommand{\aimTXT}{Целью}
\newcommand{\tasksTXT}{задачи}
% научную новизну;
\newcommand{\noveltyTXT}{Научная новизна:}
% теоретическую и практическую значимость работы;
%\newcommand{\influenceTXT}{Теоретическая и практическая значимость}
% или чаще используют просто
\newcommand{\influenceTXT}{Практическая значимость}
% методологию и методы исследования;
\newcommand{\methodsTXT}{Методология и методы исследования.}
% положения, выносимые на защиту;
\newcommand{\defpositionsTXT}{Основные положения, выносимые на~защиту:}
% степень достоверности и апробацию результатов.
\newcommand{\reliabilityTXT}{Достоверность}
\newcommand{\probationTXT}{Апробация работы.}

\newcommand{\contributionTXT}{Личный вклад.}
\newcommand{\publicationsTXT}{Публикации.}


%%% Заголовки библиографии:

% для автореферата:
\newcommand{\bibtitleauthor}{Публикации автора по теме диссертации}

% для стиля библиографии `\insertbiblioauthorgrouped`
\newcommand{\bibtitleauthorvak}{В изданиях из списка ВАК РФ}
\newcommand{\bibtitleauthorscopus}{В изданиях, входящих в международную базу цитирования Scopus}
\newcommand{\bibtitleauthorwos}{В изданиях, входящих в международную базу цитирования Web of Science}
\newcommand{\bibtitleauthorother}{В прочих изданиях}
\newcommand{\bibtitleauthorconf}{В сборниках трудов конференций}
\newcommand{\bibtitleauthorpatent}{Зарегистрированные патенты}
\newcommand{\bibtitleauthorprogram}{Зарегистрированные программы для ЭВМ}

% для стиля библиографии `\insertbiblioauthorimportant`:
\newcommand{\bibtitleauthorimportant}{Наиболее значимые \protect\MakeLowercase\bibtitleauthor}

% для списка литературы в диссертации и списка чужих работ в автореферате:
\newcommand{\bibtitlefull}{Список литературы} % (ГОСТ Р 7.0.11-2011, 4)

\newtheorem{theorem}{Теорема}
\newtheorem{lemma}{Лемма}
\newtheorem{remark}{Замечание}
\newtheorem{definition}{Определение}
       % Новые переменные, которые могут использоваться во всём проекте
%%%%%%%%%%%%%%%%%%%%%%%%%%%%%%%%%%%%%%%%%%%%%%%%%%%%%%%
%%%% Файл упрощённых настроек шаблона автореферата %%%%
%%%%%%%%%%%%%%%%%%%%%%%%%%%%%%%%%%%%%%%%%%%%%%%%%%%%%%%

%%% Инициализирование переменных, не трогать!  %%%
\newcounter{showperssign}
\newcounter{showsecrsign}
\newcounter{showopplead}
%%%%%%%%%%%%%%%%%%%%%%%%%%%%%%%%%%%%%%%%%%%%%%%%%%%%%%%

%%% Список публикаций %%%
\makeatletter
\@ifundefined{c@usefootcite}{
  \newcounter{usefootcite}
  \setcounter{usefootcite}{0} % 0 --- два списка литературы;
                              % 1 --- список публикаций автора + цитирование
                              %       других работ в сносках
}{}
\makeatother

\makeatletter
\@ifundefined{c@bibgrouped}{
  \newcounter{bibgrouped}
  \setcounter{bibgrouped}{0}  % 0 --- единый список работ автора;
                              % 1 --- сгруппированные работы автора
}{}
\makeatother

%%% Область упрощённого управления оформлением %%%

%% Управление зазором между подрисуночной подписью и основным текстом %%
\setlength{\belowcaptionskip}{10pt plus 20pt minus 2pt}


%% Подпись таблиц %%

% смещение строк подписи после первой
\newcommand{\tabindent}{0cm}

% тип форматирования таблицы
% plain --- название и текст в одной строке
% split --- название и текст в разных строках
\newcommand{\tabformat}{plain}

%%% настройки форматирования таблицы `plain'

% выравнивание по центру подписи, состоящей из одной строки
% true  --- выравнивать
% false --- не выравнивать
\newcommand{\tabsinglecenter}{false}

% выравнивание подписи таблиц
% justified   --- выравнивать как обычный текст
% centering   --- выравнивать по центру
% centerlast  --- выравнивать по центру только последнюю строку
% centerfirst --- выравнивать по центру только первую строку
% raggedleft  --- выравнивать по правому краю
% raggedright --- выравнивать по левому краю
\newcommand{\tabjust}{justified}

% Разделитель записи «Таблица #» и названия таблицы
\newcommand{\tablabelsep}{~\cyrdash\ }

%%% настройки форматирования таблицы `split'

% положение названия таблицы
% \centering   --- выравнивать по центру
% \raggedleft  --- выравнивать по правому краю
% \raggedright --- выравнивать по левому краю
\newcommand{\splitformatlabel}{\raggedleft}

% положение текста подписи
% \centering   --- выравнивать по центру
% \raggedleft  --- выравнивать по правому краю
% \raggedright --- выравнивать по левому краю
\newcommand{\splitformattext}{\raggedright}

%% Подпись рисунков %%
%Разделитель записи «Рисунок #» и названия рисунка
\newcommand{\figlabelsep}{~\cyrdash\ }  % (ГОСТ 2.105, 4.3.1)
                                        % "--- здесь не работает

%Демонстрация подписи диссертанта на автореферате
\setcounter{showperssign}{0}  % 0 --- не показывать;
                              % 1 --- показывать
%Демонстрация подписи учёного секретаря на автореферате
\setcounter{showsecrsign}{0}  % 0 --- не показывать;
                              % 1 --- показывать
%Демонстрация информации об оппонентах и ведущей организации на автореферате
\setcounter{showopplead}{1}   % 0 --- не показывать;
                              % 1 --- показывать

%%% Цвета гиперссылок %%%
% Latex color definitions: http://latexcolor.com/
\definecolor{linkcolor}{rgb}{0.9,0,0}
\definecolor{citecolor}{rgb}{0,0.6,0}
\definecolor{urlcolor}{rgb}{0,0,1}
%\definecolor{linkcolor}{rgb}{0,0,0} %black
%\definecolor{citecolor}{rgb}{0,0,0} %black
%\definecolor{urlcolor}{rgb}{0,0,0} %black
        % Упрощённые настройки шаблона

%%% Основные сведения %%%
\newcommand{\thesisAuthorLastName}{Месенёв}
\newcommand{\thesisAuthorOtherNames}{Павел Ростиславович}
\newcommand{\thesisAuthorInitials}{П.\,Р.}
\newcommand{\thesisAuthor}             % Диссертация, ФИО автора
{%
    \texorpdfstring{% \texorpdfstring takes two arguments and uses the first for (La)TeX and the second for pdf
        \thesisAuthorLastName~\thesisAuthorOtherNames% так будет отображаться на титульном листе или в тексте, где будет использоваться переменная
    }{%
        \thesisAuthorLastName, \thesisAuthorOtherNames% эта запись для свойств pdf-файла. В таком виде, если pdf будет обработан программами для сбора библиографических сведений, будет правильно представлена фамилия.
    }
}
\newcommand{\thesisAuthorShort}        % Диссертация, ФИО автора инициалами
{\thesisAuthorInitials~\thesisAuthorLastName}
%\newcommand{\thesisUdk}                % Диссертация, УДК
%{\fixme{xxx.xxx}}
\newcommand{\thesisTitle}              % Диссертация, название
{Оптимизационные методы решения обратных задач сложного теплообмена}
\newcommand{\thesisSpecialtyNumber}    % Диссертация, специальность, номер
{1.2.2}
\newcommand{\thesisSpecialtyTitle}     % Диссертация, специальность, название (название взято с сайта ВАК для примера)
{Математическое моделирование, численные методы и комплексы программ}
%% \newcommand{\thesisSpecialtyTwoNumber} % Диссертация, вторая специальность, номер
%% {\fixme{XX.XX.XX}}
%% \newcommand{\thesisSpecialtyTwoTitle}  % Диссертация, вторая специальность, название
%% {\fixme{Теория и~методика физического воспитания, спортивной тренировки,
%% оздоровительной и~адаптивной физической культуры}}
\newcommand{\thesisDegree}             % Диссертация, ученая степень
{кандидата физико-математических наук}
\newcommand{\thesisDegreeShort}        % Диссертация, ученая степень, краткая запись
{канд. физ.-мат. наук}
\newcommand{\thesisCity}               % Диссертация, город написания диссертации
{Владивосток}
\newcommand{\thesisYear}               % Диссертация, год написания диссертации
{\the\year}
\newcommand{\thesisOrganization}       % Диссертация, организация
{Федеральное государственное автономное образовательное учреждение
высшего образования «Дальневосточный федеральный университет»}
\newcommand{\thesisOrganizationShort}  % Диссертация, краткое название организации для доклада
{ДВФУ}

\newcommand{\thesisInOrganization}     % Диссертация, организация в предложном падеже: Работа выполнена в ...
{федеральном государственном автономном образовательном учреждении
высшего образования «Дальневосточный федеральный университет»}

%% \newcommand{\supervisorDead}{}           % Рисовать рамку вокруг фамилии
\newcommand{\supervisorFio}              % Научный руководитель, ФИО
{Чеботарев Александр Юрьевич}
\newcommand{\supervisorRegalia}          % Научный руководитель, регалии
{доктор физико-математических наук, профессор}
\newcommand{\supervisorFioShort}         % Научный руководитель, ФИО
{А.\,Ю.~Чеботарев}
\newcommand{\supervisorRegaliaShort}     % Научный руководитель, регалии
{д. ф.-м. н., проф.}

%% \newcommand{\supervisorTwoDead}{}        % Рисовать рамку вокруг фамилии
%% \newcommand{\supervisorTwoFio}           % Второй научный руководитель, ФИО
%% {\fixme{Фамилия Имя Отчество}}
%% \newcommand{\supervisorTwoRegalia}       % Второй научный руководитель, регалии
%% {\fixme{уч. степень, уч. звание}}
%% \newcommand{\supervisorTwoFioShort}      % Второй научный руководитель, ФИО
%% {\fixme{И.\,О.~Фамилия}}
%% \newcommand{\supervisorTwoRegaliaShort}  % Второй научный руководитель, регалии
%% {\fixme{уч.~ст.,~уч.~зв.}}

\newcommand{\opponentOneFio}           % Оппонент 1, ФИО
{Шишленин Максим Александрович}
\newcommand{\opponentOneRegalia}       % Оппонент 1, регалии
{доктор физико-математических наук, профессор РАН}
\newcommand{\opponentOneJobPlace}      % Оппонент 1, место работы
{ФГУБН Институт математики им. С. Л. Соболева СО РАН}
\newcommand{\opponentOneJobPost}       % Оппонент 1, должность
{заместитель директора по науке}

\newcommand{\opponentTwoFio}           % Оппонент 2, ФИО
{Максимова Надежда Николаевна}
\newcommand{\opponentTwoRegalia}       % Оппонент 2, регалии
{кандидат физико-математических наук, доцент}
\newcommand{\opponentTwoJobPlace}      % Оппонент 2, место работы
{ФГБОУ ВО Амурский государственный университет}
\newcommand{\opponentTwoJobPost}       % Оппонент 2, должность
{зав. кафедрой математического анализа и моделирования}

%% \newcommand{\opponentThreeFio}         % Оппонент 3, ФИО
%% {\fixme{Фамилия Имя Отчество}}
%% \newcommand{\opponentThreeRegalia}     % Оппонент 3, регалии
%% {\fixme{кандидат физико-математических наук}}
%% \newcommand{\opponentThreeJobPlace}    % Оппонент 3, место работы
%% {\fixme{Основное место работы c длинным длинным длинным длинным названием}}
%% \newcommand{\opponentThreeJobPost}     % Оппонент 3, должность
%% {\fixme{старший научный сотрудник}}

\newcommand{\leadingOrganizationTitle} % Ведущая организация, дополнительные строки. Удалить, чтобы не отображать в автореферате
{ФГБУН Тихоокеанский океанологический институт им. В.И. Ильичева ДВО РАН}

\newcommand{\defenseDate}              % Защита, дата
{\fixme{DD mmmmmmmm YYYY~г.~в~XX часов}}
\newcommand{\defenseCouncilNumber}     % Защита, номер диссертационного совета
{\fixme{Д\,123.456.78}}
\newcommand{\defenseCouncilTitle}      % Защита, учреждение диссертационного совета
{\fixme{Название учреждения}}
\newcommand{\defenseCouncilAddress}    % Защита, адрес учреждение диссертационного совета
{\fixme{Адрес}}
\newcommand{\defenseCouncilPhone}      % Телефон для справок
{\fixme{+7~(0000)~00-00-00}}

\newcommand{\defenseSecretaryFio}      % Секретарь диссертационного совета, ФИО
{\fixme{Фамилия Имя Отчество}}
\newcommand{\defenseSecretaryRegalia}  % Секретарь диссертационного совета, регалии
{\fixme{д-р~физ.-мат. наук}}            % Для сокращений есть ГОСТы, например: ГОСТ Р 7.0.12-2011 + http://base.garant.ru/179724/#block_30000

\newcommand{\synopsisLibrary}          % Автореферат, название библиотеки
{\fixme{Название библиотеки}}
\newcommand{\synopsisDate}             % Автореферат, дата рассылки
{\fixme{DD mmmmmmmm}\the\year~года}

% To avoid conflict with beamer class use \providecommand
\providecommand{\keywords}%            % Ключевые слова для метаданных PDF диссертации и автореферата
{}
           % Основные сведения
\input{common/fonts}          % Определение шрифтов (частичное)
% Общие стили оформления.
% Возможные варианты значений ищите в описании библиотеки beamer
\usetheme{Pittsburgh}
\usecolortheme{whale}

% \usetheme[secheader]{Boadilla}
% \usecolortheme{seahorse}

% Размер полей слайдов
\setbeamersize{text margin left=1cm,%
               text margin right=1cm}

% выключение кнопок навигации
\beamertemplatenavigationsymbolsempty

% Размеры шрифтов
\setbeamerfont{title}{size=\large}
\setbeamerfont{subtitle}{size=\small}
\setbeamerfont{author}{size=\normalsize}
\setbeamerfont{institute}{size=\small}
\setbeamerfont{date}{size=\normalsize}
\setbeamerfont{bibliography item}{size=\small}
\setbeamerfont{bibliography entry author}{size=\small}
\setbeamerfont{bibliography entry title}{size=\small}
\setbeamerfont{bibliography entry location}{size=\small}
\setbeamerfont{bibliography entry note}{size=\small}
\setbeamercolor{math text}{fg=black!15!blue}
% Аналогично можно настроить и другие размеры.
% Названия классов элементов можно найти здесь
% http://www.cpt.univ-mrs.fr/~masson/latex/Beamer-appearance-cheat-sheet.pdf

% Цвет элементов
\setbeamercolor{footline}{fg=blue}
\setbeamercolor{bibliography item}{fg=black}
\setbeamercolor{bibliography entry author}{fg=black}
\setbeamercolor{bibliography entry title}{fg=black}
\setbeamercolor{bibliography entry location}{fg=black}
\setbeamercolor{bibliography entry note}{fg=black}
% Аналогично можно настроить и другие цвета.
% Названия классов элементов можно найти здесь
% http://www.cpt.univ-mrs.fr/~masson/latex/Beamer-appearance-cheat-sheet.pdf

% Нумеровать список статей
% https://tex.stackexchange.com/a/419506/104425
\setbeamertemplate{bibliography item}{\insertbiblabel}
% или убрать номера
% \setbeamertemplate{bibliography item}{}

% Использовать шрифт с засечками для формул
% https://tex.stackexchange.com/a/34267/104425
\usefonttheme[onlymath]{serif}

% https://tex.stackexchange.com/a/291545/104425
\makeatletter
\def\beamer@framenotesbegin{% at beginning of slide
    \usebeamercolor[fg]{normal text}
    \gdef\beamer@noteitems{}%
    \gdef\beamer@notes{}%
}
\makeatother

% footer презентации
\setbeamertemplate{footline}{
    \leavevmode%
    \hbox{%
        \begin{beamercolorbox}[wd=.333333\paperwidth,ht=2.25ex,dp=1ex,center]{}%
            % И. О. Фамилия, Организация кратко
            \thesisAuthorShort, \thesisOrganizationShort
        \end{beamercolorbox}%
        \begin{beamercolorbox}[wd=.333333\paperwidth,ht=2.25ex,dp=1ex,center]{}%
            % Город, 20XX
            \thesisCity, \thesisYear
        \end{beamercolorbox}%
        \begin{beamercolorbox}[wd=.333333\paperwidth,ht=2.25ex,dp=1ex,right]{}%
            Стр. \insertframenumber{} из \inserttotalframenumber \hspace*{2ex}
        \end{beamercolorbox}}%
    \vskip0pt%
}

% вывод на экран заметок к презентации
\ifnumequal{\value{presnotes}}{0}{}{
    \setbeameroption{show notes}
    \ifnumequal{\value{presnotes}}{2}{
        \setbeameroption{show notes on second screen=\presposition}
    }{}
}
         % Стили общие для диссертации и автореферата
\input{Synopsis/synstyles}    % Стили для автореферата
% для вертикального центрирования ячеек в tabulary
\def\zz{\ifx\[$\else\aftergroup\zzz\fi}
%$ \] % <-- чиним подсветку синтаксиса в некоторых редакторах
\def\zzz{\setbox0\lastbox
\dimen0\dimexpr\extrarowheight + \ht0-\dp0\relax
\setbox0\hbox{\raise-.5\dimen0\box0}%
\ht0=\dimexpr\ht0+\extrarowheight\relax
\dp0=\dimexpr\dp0+\extrarowheight\relax
\box0
}

\lstdefinelanguage{Renhanced}%
{keywords={abbreviate,abline,abs,acos,acosh,action,add1,add,%
        aggregate,alias,Alias,alist,all,anova,any,aov,aperm,append,apply,%
        approx,approxfun,apropos,Arg,args,array,arrows,as,asin,asinh,%
        atan,atan2,atanh,attach,attr,attributes,autoload,autoloader,ave,%
        axis,backsolve,barplot,basename,besselI,besselJ,besselK,besselY,%
        beta,binomial,body,box,boxplot,break,browser,bug,builtins,bxp,by,%
        c,C,call,Call,case,cat,category,cbind,ceiling,character,char,%
        charmatch,check,chol,chol2inv,choose,chull,class,close,cm,codes,%
        coef,coefficients,co,col,colnames,colors,colours,commandArgs,%
        comment,complete,complex,conflicts,Conj,contents,contour,%
        contrasts,contr,control,helmert,contrib,convolve,cooks,coords,%
        distance,coplot,cor,cos,cosh,count,fields,cov,covratio,wt,CRAN,%
        create,crossprod,cummax,cummin,cumprod,cumsum,curve,cut,cycle,D,%
        data,dataentry,date,dbeta,dbinom,dcauchy,dchisq,de,debug,%
        debugger,Defunct,default,delay,delete,deltat,demo,de,density,%
        deparse,dependencies,Deprecated,deriv,description,detach,%
        dev2bitmap,dev,cur,deviance,off,prev,,dexp,df,dfbetas,dffits,%
        dgamma,dgeom,dget,dhyper,diag,diff,digamma,dim,dimnames,dir,%
        dirname,dlnorm,dlogis,dnbinom,dnchisq,dnorm,do,dotplot,double,%
        download,dpois,dput,drop,drop1,dsignrank,dt,dummy,dump,dunif,%
        duplicated,dweibull,dwilcox,dyn,edit,eff,effects,eigen,else,%
        emacs,end,environment,env,erase,eval,equal,evalq,example,exists,%
        exit,exp,expand,expression,External,extract,extractAIC,factor,%
        fail,family,fft,file,filled,find,fitted,fivenum,fix,floor,for,%
        For,formals,format,formatC,formula,Fortran,forwardsolve,frame,%
        frequency,ftable,ftable2table,function,gamma,Gamma,gammaCody,%
        gaussian,gc,gcinfo,gctorture,get,getenv,geterrmessage,getOption,%
        getwd,gl,glm,globalenv,gnome,GNOME,graphics,gray,grep,grey,grid,%
        gsub,hasTsp,hat,heat,help,hist,home,hsv,httpclient,I,identify,if,%
        ifelse,Im,image,\%in\%,index,influence,measures,inherits,install,%
        installed,integer,interaction,interactive,Internal,intersect,%
        inverse,invisible,IQR,is,jitter,kappa,kronecker,labels,lapply,%
        layout,lbeta,lchoose,lcm,legend,length,levels,lgamma,library,%
        licence,license,lines,list,lm,load,local,locator,log,log10,log1p,%
        log2,logical,loglin,lower,lowess,ls,lsfit,lsf,ls,machine,Machine,%
        mad,mahalanobis,make,link,margin,match,Math,matlines,mat,matplot,%
        matpoints,matrix,max,mean,median,memory,menu,merge,methods,min,%
        missing,Mod,mode,model,response,mosaicplot,mtext,mvfft,na,nan,%
        names,omit,nargs,nchar,ncol,NCOL,new,next,NextMethod,nextn,%
        nlevels,nlm,noquote,NotYetImplemented,NotYetUsed,nrow,NROW,null,%
        numeric,\%o\%,objects,offset,old,on,Ops,optim,optimise,optimize,%
        options,or,order,ordered,outer,package,packages,page,pairlist,%
        pairs,palette,panel,par,parent,parse,paste,path,pbeta,pbinom,%
        pcauchy,pchisq,pentagamma,persp,pexp,pf,pgamma,pgeom,phyper,pico,%
        pictex,piechart,Platform,plnorm,plogis,plot,pmatch,pmax,pmin,%
        pnbinom,pnchisq,pnorm,points,poisson,poly,polygon,polyroot,pos,%
        postscript,power,ppoints,ppois,predict,preplot,pretty,Primitive,%
        print,prmatrix,proc,prod,profile,proj,prompt,prop,provide,%
        psignrank,ps,pt,ptukey,punif,pweibull,pwilcox,q,qbeta,qbinom,%
        qcauchy,qchisq,qexp,qf,qgamma,qgeom,qhyper,qlnorm,qlogis,qnbinom,%
        qnchisq,qnorm,qpois,qqline,qqnorm,qqplot,qr,Q,qty,qy,qsignrank,%
        qt,qtukey,quantile,quasi,quit,qunif,quote,qweibull,qwilcox,%
        rainbow,range,rank,rbeta,rbind,rbinom,rcauchy,rchisq,Re,read,csv,%
        csv2,fwf,readline,socket,real,Recall,rect,reformulate,regexpr,%
        relevel,remove,rep,repeat,replace,replications,report,require,%
        resid,residuals,restart,return,rev,rexp,rf,rgamma,rgb,rgeom,R,%
        rhyper,rle,rlnorm,rlogis,rm,rnbinom,RNGkind,rnorm,round,row,%
        rownames,rowsum,rpois,rsignrank,rstandard,rstudent,rt,rug,runif,%
        rweibull,rwilcox,sample,sapply,save,scale,scan,scan,screen,sd,se,%
        search,searchpaths,segments,seq,sequence,setdiff,setequal,set,%
        setwd,show,sign,signif,sin,single,sinh,sink,solve,sort,source,%
        spline,splinefun,split,sqrt,stars,start,stat,stem,step,stop,%
        storage,strstrheight,stripplot,strsplit,structure,strwidth,sub,%
        subset,substitute,substr,substring,sum,summary,sunflowerplot,svd,%
        sweep,switch,symbol,symbols,symnum,sys,status,system,t,table,%
        tabulate,tan,tanh,tapply,tempfile,terms,terrain,tetragamma,text,%
        time,title,topo,trace,traceback,transform,tri,trigamma,trunc,try,%
        ts,tsp,typeof,unclass,undebug,undoc,union,unique,uniroot,unix,%
        unlink,unlist,unname,untrace,update,upper,url,UseMethod,var,%
        variable,vector,Version,vi,warning,warnings,weighted,weights,%
        which,while,window,write,\%x\%,x11,X11,xedit,xemacs,xinch,xor,%
        xpdrows,xy,xyinch,yinch,zapsmall,zip},%
    otherkeywords={!,!=,~,$,*,\%,\&,\%/\%,\%*\%,\%\%,<-,<<-},%$
    alsoother={._$},%$
    sensitive,%
    morecomment=[l]\#,%
    morestring=[d]",%
    morestring=[d]'% 2001 Robert Denham
}%

%решаем проблему с кириллицей в комментариях (в pdflatex) https://tex.stackexchange.com/a/103712
\lstset{extendedchars=true,keepspaces=true,literate={Ö}{{\"O}}1
    {Ä}{{\"A}}1
    {Ü}{{\"U}}1
    {ß}{{\ss}}1
    {ü}{{\"u}}1
    {ä}{{\"a}}1
    {ö}{{\"o}}1
    {~}{{\textasciitilde}}1
    {а}{{\selectfont\char224}}1
    {б}{{\selectfont\char225}}1
    {в}{{\selectfont\char226}}1
    {г}{{\selectfont\char227}}1
    {д}{{\selectfont\char228}}1
    {е}{{\selectfont\char229}}1
    {ё}{{\"e}}1
    {ж}{{\selectfont\char230}}1
    {з}{{\selectfont\char231}}1
    {и}{{\selectfont\char232}}1
    {й}{{\selectfont\char233}}1
    {к}{{\selectfont\char234}}1
    {л}{{\selectfont\char235}}1
    {м}{{\selectfont\char236}}1
    {н}{{\selectfont\char237}}1
    {о}{{\selectfont\char238}}1
    {п}{{\selectfont\char239}}1
    {р}{{\selectfont\char240}}1
    {с}{{\selectfont\char241}}1
    {т}{{\selectfont\char242}}1
    {у}{{\selectfont\char243}}1
    {ф}{{\selectfont\char244}}1
    {х}{{\selectfont\char245}}1
    {ц}{{\selectfont\char246}}1
    {ч}{{\selectfont\char247}}1
    {ш}{{\selectfont\char248}}1
    {щ}{{\selectfont\char249}}1
    {ъ}{{\selectfont\char250}}1
    {ы}{{\selectfont\char251}}1
    {ь}{{\selectfont\char252}}1
    {э}{{\selectfont\char253}}1
    {ю}{{\selectfont\char254}}1
    {я}{{\selectfont\char255}}1
    {А}{{\selectfont\char192}}1
    {Б}{{\selectfont\char193}}1
    {В}{{\selectfont\char194}}1
    {Г}{{\selectfont\char195}}1
    {Д}{{\selectfont\char196}}1
    {Е}{{\selectfont\char197}}1
    {Ё}{{\"E}}1
    {Ж}{{\selectfont\char198}}1
    {З}{{\selectfont\char199}}1
    {И}{{\selectfont\char200}}1
    {Й}{{\selectfont\char201}}1
    {К}{{\selectfont\char202}}1
    {Л}{{\selectfont\char203}}1
    {М}{{\selectfont\char204}}1
    {Н}{{\selectfont\char205}}1
    {О}{{\selectfont\char206}}1
    {П}{{\selectfont\char207}}1
    {Р}{{\selectfont\char208}}1
    {С}{{\selectfont\char209}}1
    {Т}{{\selectfont\char210}}1
    {У}{{\selectfont\char211}}1
    {Ф}{{\selectfont\char212}}1
    {Х}{{\selectfont\char213}}1
    {Ц}{{\selectfont\char214}}1
    {Ч}{{\selectfont\char215}}1
    {Ш}{{\selectfont\char216}}1
    {Щ}{{\selectfont\char217}}1
    {Ъ}{{\selectfont\char218}}1
    {Ы}{{\selectfont\char219}}1
    {Ь}{{\selectfont\char220}}1
    {Э}{{\selectfont\char221}}1
    {Ю}{{\selectfont\char222}}1
    {Я}{{\selectfont\char223}}1
    {і}{{\selectfont\char105}}1
    {ї}{{\selectfont\char168}}1
    {є}{{\selectfont\char185}}1
    {ґ}{{\selectfont\char160}}1
    {І}{{\selectfont\char73}}1
    {Ї}{{\selectfont\char136}}1
    {Є}{{\selectfont\char153}}1
    {Ґ}{{\selectfont\char128}}1
}

% Ширина текста минус ширина надписи 999
\newlength{\twless}
\newlength{\lmarg}
\setlength{\lmarg}{\widthof{999}}   % ширина надписи 999
\setlength{\twless}{\textwidth-\lmarg}

\lstset{ %
%    language=R,                     %  Язык указать здесь, если во всех листингах преимущественно один язык, в результате часть настроек может пойти только для этого языка
    numbers=left,                   % where to put the line-numbers
    numberstyle=\fontsize{12pt}{14pt}\selectfont\color{Gray},  % the style that is used for the line-numbers
    firstnumber=1,                  % в этой и следующей строках задаётся поведение нумерации 5, 10, 15...
    stepnumber=5,                   % the step between two line-numbers. If it's 1, each line will be numbered
    numbersep=5pt,                  % how far the line-numbers are from the code
    backgroundcolor=\color{white},  % choose the background color. You must add \usepackage{color}
    showspaces=false,               % show spaces adding particular underscores
    showstringspaces=false,         % underline spaces within strings
    showtabs=false,                 % show tabs within strings adding particular underscores
    frame=leftline,                 % adds a frame of different types around the code
    rulecolor=\color{black},        % if not set, the frame-color may be changed on line-breaks within not-black text (e.g. commens (green here))
    tabsize=2,                      % sets default tabsize to 2 spaces
    captionpos=t,                   % sets the caption-position to top
    breaklines=true,                % sets automatic line breaking
    breakatwhitespace=false,        % sets if automatic breaks should only happen at whitespace
%    title=\lstname,                 % show the filename of files included with \lstinputlisting;
    % also try caption instead of title
    basicstyle=\fontsize{12pt}{14pt}\selectfont\ttfamily,% the size of the fonts that are used for the code
%    keywordstyle=\color{blue},      % keyword style
    commentstyle=\color{ForestGreen}\emph,% comment style
    stringstyle=\color{Mahogany},   % string literal style
    escapeinside={\%*}{*)},         % if you want to add a comment within your code
    morekeywords={*,...},           % if you want to add more keywords to the set
    inputencoding=utf8,             % кодировка кода
    xleftmargin={\lmarg},           % Чтобы весь код и полоска с номерами строк была смещена влево, так чтобы цифры не вылезали за пределы текста слева
}

%http://tex.stackexchange.com/questions/26872/smaller-frame-with-listings
% Окружение, чтобы листинг был компактнее обведен рамкой, если она задается, а не на всю ширину текста
\makeatletter
\newenvironment{SmallListing}[1][]
{\lstset{#1}\VerbatimEnvironment\begin{VerbatimOut}{VerbEnv.tmp}}
{\end{VerbatimOut}\settowidth\@tempdima{%
        \lstinputlisting{VerbEnv.tmp}}
    \minipage{\@tempdima}\lstinputlisting{VerbEnv.tmp}\endminipage}
\makeatother

\DefineVerbatimEnvironment% с шрифтом 12 пт
{Verb}{Verbatim}
{fontsize=\fontsize{12pt}{14pt}\selectfont}

\newfloat[chapter]{ListingEnv}{lol}{Листинг}

\renewcommand{\lstlistingname}{Листинг}

%Общие счётчики окружений листингов
%http://tex.stackexchange.com/questions/145546/how-to-make-figure-and-listing-share-their-counter
% Если смешивать плавающие и не плавающие окружения, то могут быть проблемы с нумерацией
\makeatletter
\AfterEndPreamble{% https://tex.stackexchange.com/a/252682
    \let\c@ListingEnv\relax % drop existing counter "ListingEnv"
    \newaliascnt{ListingEnv}{lstlisting} % команда требует пакет aliascnt
    \let\ftype@lstlisting\ftype@ListingEnv % give the floats the same precedence
}
\makeatother

% значок С++ — используйте команду \cpp
\newcommand{\cpp}{%
    C\nolinebreak\hspace{-.05em}%
    \raisebox{.2ex}{+}\nolinebreak\hspace{-.10em}%
    \raisebox{.2ex}{+}%
}

%%% Ради примера во второй главе
\let\originalepsilon\epsilon
\let\originalphi\phi
\let\originalkappa\kappa
\let\originalle\le
\let\originalleq\leq
\let\originalge\ge
\let\originalgeq\geq
\let\originalemptyset\emptyset
\let\originaltan\tan
\let\originalcot\cot
\let\originalcsc\csc

%%% Русская традиция начертания математических знаков
\renewcommand{\le}{\ensuremath{\leqslant}}
\renewcommand{\leq}{\ensuremath{\leqslant}}
\renewcommand{\ge}{\ensuremath{\geqslant}}
\renewcommand{\geq}{\ensuremath{\geqslant}}
\renewcommand{\emptyset}{\varnothing}

%%% Русская традиция начертания математических функций (на случай копирования из зарубежных источников)
\renewcommand{\tan}{\operatorname{tg}}
\renewcommand{\cot}{\operatorname{ctg}}
\renewcommand{\csc}{\operatorname{cosec}}

%%% Русская традиция начертания греческих букв (греческие буквы вертикальные, через пакет upgreek)
\renewcommand{\epsilon}{\ensuremath{\upvarepsilon}}   %  русская традиция записи
\renewcommand{\phi}{\ensuremath{\upvarphi}}
%\renewcommand{\kappa}{\ensuremath{\varkappa}}
\renewcommand{\alpha}{\upalpha}
\renewcommand{\beta}{\upbeta}
\renewcommand{\gamma}{\upgamma}
\renewcommand{\delta}{\updelta}
\renewcommand{\varepsilon}{\upvarepsilon}
\renewcommand{\zeta}{\upzeta}
\renewcommand{\eta}{\upeta}
\renewcommand{\theta}{\uptheta}
\renewcommand{\vartheta}{\upvartheta}
\renewcommand{\iota}{\upiota}
\renewcommand{\kappa}{\upkappa}
\renewcommand{\lambda}{\uplambda}
\renewcommand{\mu}{\upmu}
\renewcommand{\nu}{\upnu}
\renewcommand{\xi}{\upxi}
\renewcommand{\pi}{\uppi}
\renewcommand{\varpi}{\upvarpi}
\renewcommand{\rho}{\uprho}
%\renewcommand{\varrho}{\upvarrho}
\renewcommand{\sigma}{\upsigma}
%\renewcommand{\varsigma}{\upvarsigma}
\renewcommand{\tau}{\uptau}
\renewcommand{\upsilon}{\upupsilon}
\renewcommand{\varphi}{\upvarphi}
\renewcommand{\chi}{\upchi}
\renewcommand{\psi}{\uppsi}
\renewcommand{\omega}{\upomega}
   % Стили для специфических пользовательских задач

%%%% Библиография. Выбор движка для реализации %%%
%\ifnumequal{\value{bibliosel}}{0}{%
%    \input{biblio/predefined} % Встроенная реализация с загрузкой файла через движок bibtex8
%}{
    %%% Реализация библиографии пакетами biblatex и biblatex-gost с использованием движка biber %%%

\usepackage{csquotes} % biblatex рекомендует его подключать. Пакет для оформления сложных блоков цитирования.
%%% Загрузка пакета с основными настройками %%%
\makeatletter
\ifnumequal{\value{draft}}{0}{% Чистовик
\usepackage[%
backend=biber,% движок
bibencoding=utf8,% кодировка bib файла
sorting=none,% настройка сортировки списка литературы
style=numeric-comp,% стиль цитирования и библиографии (по ГОСТ)
language=autobib,% получение языка из babel/polyglossia, default: autobib % если ставить autocite или auto, то цитаты в тексте с указанием страницы, получат указание страницы на языке оригинала
autolang=other,% многоязычная библиография
clearlang=true,% внутренний сброс поля language, если он совпадает с языком из babel/polyglossia
defernumbers=true,% нумерация проставляется после двух компиляций, зато позволяет выцеплять библиографию по ключевым словам и нумеровать не из большего списка
sortcites=true,% сортировать номера затекстовых ссылок при цитировании (если в квадратных скобках несколько ссылок, то отображаться будут отсортированно, а не абы как)
doi=false,% Показывать или нет ссылки на DOI
isbn=false,% Показывать или нет ISBN, ISSN, ISRN
]{biblatex}[2016/09/17]
\ltx@iffilelater{biblatex-gost.def}{2017/05/03}%
{\toggletrue{bbx:gostbibliography}%
\renewcommand*{\revsdnamepunct}{\addcomma}}{}
}{%Черновик
\usepackage[%
backend=biber,% движок
bibencoding=utf8,% кодировка bib файла
sorting=none,% настройка сортировки списка литературы
style=numeric-comp,% стиль цитирования и библиографии (по ГОСТ)
% defernumbers=true, % откомментируйте, если требуется правильная нумерация ссылок на литературу в режиме черновика. Замедляет сборку
]{biblatex}[2016/09/17]%
}
\makeatother

\providebool{blxmc} % biblatex version needs and has MakeCapital workaround
\boolfalse{blxmc} % setting our new boolean flag to default false
\ifxetexorluatex
\else
% Исправление случая неподдержки знака номера в pdflatex
    \DefineBibliographyStrings{russian}{number={\textnumero}}

% Исправление случая отсутствия прописных букв в некоторых случаях
% https://github.com/plk/biblatex/issues/960#issuecomment-596658282
    \ifdefmacro{\ExplSyntaxOn}{}{\usepackage{expl3}}
    \makeatletter
    \ltx@ifpackagelater{biblatex}{2020/02/23}{
    % Assuming this version of biblatex defines MakeCapital correctly
    }{
        \ltx@ifpackagelater{biblatex}{2019/12/01}{
            % Assuming this version of biblatex defines MakeCapital incorrectly
            \usepackage{expl3}[2020/02/25]
            \@ifpackagelater{expl3}{2020/02/25}{
                \booltrue{blxmc} % setting our new boolean flag to true
            }{}
        }{}
    }
    \makeatother
    \ifblxmc
        \typeout{Assuming this version of biblatex defines MakeCapital
        incorrectly}
        \usepackage{xparse}
        \makeatletter
        \ExplSyntaxOn
        \NewDocumentCommand \blx@maketext@lowercase {m}
          {
            \text_lowercase:n {#1}
          }

        \NewDocumentCommand \blx@maketext@uppercase {m}
          {
            \text_uppercase:n {#1}
          }

        \RenewDocumentCommand \MakeCapital {m}
          {
            \text_titlecase_first:n {#1}
          }
        \ExplSyntaxOff

        \protected\def\blx@biblcstring#1#2#3{%
          \blx@begunit
          \blx@hyphenreset
          \blx@bibstringsimple
          \lowercase{\edef\blx@tempa{#3}}%
          \ifcsundef{#2@\blx@tempa}
            {\blx@warn@nostring\blx@tempa
             \blx@endnounit}
            {#1{\blx@maketext@lowercase{\csuse{#2@\blx@tempa}}}%
             \blx@endunit}}

        \protected\def\blx@bibucstring#1#2#3{%
          \blx@begunit
          \blx@hyphenreset
          \blx@bibstringsimple
          \lowercase{\edef\blx@tempa{#3}}%
          \ifcsundef{#2@\blx@tempa}
            {\blx@warn@nostring\blx@tempa
             \blx@endnounit}
            {#1{\blx@maketext@uppercase{\csuse{#2@\blx@tempa}}}%
             \blx@endunit}}
        \makeatother
    \fi
\fi

\ifsynopsis
\ifnumgreater{\value{usefootcite}}{0}{
    \ExecuteBibliographyOptions{autocite=footnote}
    \newbibmacro*{cite:full}{%
        \printtext[bibhypertarget]{%
            \usedriver{%
                \DeclareNameAlias{sortname}{default}%
            }{%
                \thefield{entrytype}%
            }%
        }%
        \usebibmacro{shorthandintro}%
    }
    \DeclareCiteCommand{\smartcite}[\mkbibfootnote]{%
        \usebibmacro{prenote}%
    }{%
        \usebibmacro{citeindex}%
        \usebibmacro{cite:full}%
    }{%
        \multicitedelim%
    }{%
        \usebibmacro{postnote}%
    }
}{}
\fi

%%% Подключение файлов bib %%%
\addbibresource[label=bl-external]{biblio/external.bib}
\addbibresource[label=bl-grenkin]{biblio/grenkin.bib}
\addbibresource[label=bl-evla]{biblio/evla.bib}
\addbibresource[label=bl-author]{biblio/author.bib}
\addbibresource[label=bl-2_4]{biblio/2_4.bib}
\addbibresource[label=bl-registered]{biblio/registered.bib}

%http://tex.stackexchange.com/a/141831/79756
%There is a way to automatically map the language field to the langid field. The following lines in the preamble should be enough to do that.
%This command will copy the language field into the langid field and will then delete the contents of the language field. The language field will only be deleted if it was successfully copied into the langid field.
\DeclareSourcemap{ %модификация bib файла перед тем, как им займётся biblatex
    \maps{
        \map{% перекидываем значения полей language в поля langid, которыми пользуется biblatex
            \step[fieldsource=language, fieldset=langid, origfieldval, final]
            \step[fieldset=language, null]
        }
        \map{% перекидываем значения полей numpages в поля pagetotal, которыми пользуется biblatex
            \step[fieldsource=numpages, fieldset=pagetotal, origfieldval, final]
            \step[fieldset=numpages, null]
        }
        \map{% перекидываем значения полей pagestotal в поля pagetotal, которыми пользуется biblatex
            \step[fieldsource=pagestotal, fieldset=pagetotal, origfieldval, final]
            \step[fieldset=pagestotal, null]
        }
        \map[overwrite]{% перекидываем значения полей shortjournal, если они есть, в поля journal, которыми пользуется biblatex
            \step[fieldsource=shortjournal, final]
            \step[fieldset=journal, origfieldval]
            \step[fieldset=shortjournal, null]
        }
        \map[overwrite]{% перекидываем значения полей shortbooktitle, если они есть, в поля booktitle, которыми пользуется biblatex
            \step[fieldsource=shortbooktitle, final]
            \step[fieldset=booktitle, origfieldval]
            \step[fieldset=shortbooktitle, null]
        }
        \map{% если в поле medium написано "Электронный ресурс", то устанавливаем поле media, которым пользуется biblatex, в значение eresource.
            \step[fieldsource=medium,
            match=\regexp{Электронный\s+ресурс},
            final]
            \step[fieldset=media, fieldvalue=eresource]
            \step[fieldset=medium, null]
        }
        \map[overwrite]{% стираем значения всех полей issn
            \step[fieldset=issn, null]
        }
        \map[overwrite]{% стираем значения всех полей abstract, поскольку ими не пользуемся, а там бывают "неприятные" латеху символы
            \step[fieldsource=abstract]
            \step[fieldset=abstract,null]
        }
        \map[overwrite]{ % переделка формата записи даты
            \step[fieldsource=urldate,
            match=\regexp{([0-9]{2})\.([0-9]{2})\.([0-9]{4})},
            replace={$3-$2-$1$4}, % $4 вставлен исключительно ради нормальной работы программ подсветки синтаксиса, которые некорректно обрабатывают $ в таких конструкциях
            final]
        }
        \map[overwrite]{ % стираем ключевые слова
            \step[fieldsource=keywords]
            \step[fieldset=keywords,null]
        }
        % реализация foreach различается для biblatex v3.12 и v3.13.
        % Для версии v3.13 эта конструкция заменяет последующие 7 структур map
        % \map[overwrite,foreach={authorvak,authorscopus,authorwos,authorconf,authorother,authorparent,authorprogram}]{ % записываем информацию о типе публикации в ключевые слова
        %     \step[fieldsource=$MAPLOOP,final=true]
        %     \step[fieldset=keywords,fieldvalue={,biblio$MAPLOOP},append=true]
        % }
        \map[overwrite]{ % записываем информацию о типе публикации в ключевые слова
            \step[fieldsource=authorvak,final=true]
            \step[fieldset=keywords,fieldvalue={,biblioauthorvak},append=true]
        }
        \map[overwrite]{ % записываем информацию о типе публикации в ключевые слова
            \step[fieldsource=authorscopus,final=true]
            \step[fieldset=keywords,fieldvalue={,biblioauthorscopus},append=true]
        }
        \map[overwrite]{ % записываем информацию о типе публикации в ключевые слова
            \step[fieldsource=authorwos,final=true]
            \step[fieldset=keywords,fieldvalue={,biblioauthorwos},append=true]
        }
        \map[overwrite]{ % записываем информацию о типе публикации в ключевые слова
            \step[fieldsource=authorconf,final=true]
            \step[fieldset=keywords,fieldvalue={,biblioauthorconf},append=true]
        }
        \map[overwrite]{ % записываем информацию о типе публикации в ключевые слова
            \step[fieldsource=authorother,final=true]
            \step[fieldset=keywords,fieldvalue={,biblioauthorother},append=true]
        }
        \map[overwrite]{ % записываем информацию о типе публикации в ключевые слова
            \step[fieldsource=authorpatent,final=true]
            \step[fieldset=keywords,fieldvalue={,biblioauthorpatent},append=true]
        }
        \map[overwrite]{ % записываем информацию о типе публикации в ключевые слова
            \step[fieldsource=authorprogram,final=true]
            \step[fieldset=keywords,fieldvalue={,biblioauthorprogram},append=true]
        }
        \map[overwrite]{ % добавляем ключевые слова, чтобы различать источники
            \perdatasource{biblio/external.bib}
            \step[fieldset=keywords, fieldvalue={,biblioexternal},append=true]
        }
        \map[overwrite]{ % добавляем ключевые слова, чтобы различать источники
            \perdatasource{biblio/author.bib}
            \step[fieldset=keywords, fieldvalue={,biblioauthor},append=true]
        }
        \map[overwrite]{ % добавляем ключевые слова, чтобы различать источники
            \perdatasource{biblio/registered.bib}
            \step[fieldset=keywords, fieldvalue={,biblioregistered},append=true]
        }
        \map[overwrite]{ % добавляем ключевые слова, чтобы различать источники
            \step[fieldset=keywords, fieldvalue={,bibliofull},append=true]
        }
%        \map[overwrite]{% стираем значения всех полей series
%            \step[fieldset=series, null]
%        }
        \map[overwrite]{% перекидываем значения полей howpublished в поля organization для типа online
            \step[typesource=online, typetarget=online, final]
            \step[fieldsource=howpublished, fieldset=organization, origfieldval]
            \step[fieldset=howpublished, null]
        }
    }
}

\ifnumequal{\value{mediadisplay}}{1}{
    \DeclareSourcemap{
        \maps{%
            \map{% использование media=text по умолчанию
                \step[fieldset=media, fieldvalue=text]
            }
        }
    }
}{}
\ifnumequal{\value{mediadisplay}}{2}{
    \DeclareSourcemap{
        \maps{%
            \map[overwrite]{% удаление всех записей media
                \step[fieldset=media, null]
            }
        }
    }
}{}
\ifnumequal{\value{mediadisplay}}{3}{
    \DeclareSourcemap{
        \maps{
            \map[overwrite]{% стираем значения всех полей media=text
                \step[fieldsource=media,match={text},final]
                \step[fieldset=media, null]
            }
        }
    }
}{}
\ifnumequal{\value{mediadisplay}}{4}{
    \DeclareSourcemap{
        \maps{
            \map[overwrite]{% стираем значения всех полей media=eresource
                \step[fieldsource=media,match={eresource},final]
                \step[fieldset=media, null]
            }
        }
    }
}{}

\ifsynopsis
\else
\DeclareSourcemap{ %модификация bib файла перед тем, как им займётся biblatex
    \maps{
        \map[overwrite]{% стираем значения всех полей addendum
            \perdatasource{biblio/author.bib}
            \step[fieldset=addendum, null] %чтобы избавиться от информации об объёме авторских статей, в отличие от автореферата
        }
    }
}
\fi

\ifpresentation
% удаляем лишние поля в списке литературы презентации
% их названия можно узнать в файле presentation.bbl
\DeclareSourcemap{
    \maps{
    \map[overwrite,foreach={%
        % {{{ Список лишних полей в презентации
        address,%
        chapter,%
        edition,%
        editor,%
        eid,%
        howpublished,%
        institution,%
        key,%
        month,%
        note,%
        number,%
        organization,%
        pages,%
        publisher,%
        school,%
        series,%
        type,%
        media,%
        url,%
        doi,%
        location,%
        volume,%
        % Список лишних полей в презентации }}}
    }]{
        \perdatasource{biblio/author.bib}
        \step[fieldset=$MAPLOOP,null]
    }
    }
}
\fi

\defbibfilter{vakscopuswos}{%
    keyword=biblioauthorvak or keyword=biblioauthorscopus or keyword=biblioauthorwos
}

\defbibfilter{scopuswos}{%
    keyword=biblioauthorscopus or keyword=biblioauthorwos
}

\defbibfilter{papersregistered}{%
    keyword=biblioauthor or keyword=biblioregistered
}

%%% Убираем неразрывные пробелы перед двоеточием и точкой с запятой %%%
%\makeatletter
%\ifnumequal{\value{draft}}{0}{% Чистовик
%    \renewcommand*{\addcolondelim}{%
%      \begingroup%
%      \def\abx@colon{%
%        \ifdim\lastkern>\z@\unkern\fi%
%        \abx@puncthook{:}\space}%
%      \addcolon%
%      \endgroup}
%
%    \renewcommand*{\addsemicolondelim}{%
%      \begingroup%
%      \def\abx@semicolon{%
%        \ifdim\lastkern>\z@\unkern\fi%
%        \abx@puncthook{;}\space}%
%      \addsemicolon%
%      \endgroup}
%}{}
%\makeatother

%%% Правка записей типа thesis, чтобы дважды не писался автор
%\ifnumequal{\value{draft}}{0}{% Чистовик
%\DeclareBibliographyDriver{thesis}{%
%  \usebibmacro{bibindex}%
%  \usebibmacro{begentry}%
%  \usebibmacro{heading}%
%  \newunit
%  \usebibmacro{author}%
%  \setunit*{\labelnamepunct}%
%  \usebibmacro{thesistitle}%
%  \setunit{\respdelim}%
%  %\printnames[last-first:full]{author}%Вот эту строчку нужно убрать, чтобы автор диссертации не дублировался
%  \newunit\newblock
%  \printlist[semicolondelim]{specdata}%
%  \newunit
%  \usebibmacro{institution+location+date}%
%  \newunit\newblock
%  \usebibmacro{chapter+pages}%
%  \newunit
%  \printfield{pagetotal}%
%  \newunit\newblock
%  \usebibmacro{doi+eprint+url+note}%
%  \newunit\newblock
%  \usebibmacro{addendum+pubstate}%
%  \setunit{\bibpagerefpunct}\newblock
%  \usebibmacro{pageref}%
%  \newunit\newblock
%  \usebibmacro{related:init}%
%  \usebibmacro{related}%
%  \usebibmacro{finentry}}
%}{}

%\newbibmacro{string+doi}[1]{% новая макрокоманда на простановку ссылки на doi
%    \iffieldundef{doi}{#1}{\href{http://dx.doi.org/\thefield{doi}}{#1}}}

%\ifnumequal{\value{draft}}{0}{% Чистовик
%\renewcommand*{\mkgostheading}[1]{\usebibmacro{string+doi}{#1}} % ссылка на doi с авторов. стоящих впереди записи
%\renewcommand*{\mkgostheading}[1]{#1} % только лишь убираем курсив с авторов
%}{}
%\DeclareFieldFormat{title}{\usebibmacro{string+doi}{#1}} % ссылка на doi с названия работы
%\DeclareFieldFormat{journaltitle}{\usebibmacro{string+doi}{#1}} % ссылка на doi с названия журнала
%%% Тире как разделитель в библиографии традиционной руской длины:
\renewcommand*{\newblockpunct}{\addperiod\addnbspace\cyrdash\space\bibsentence}
%%% Убрать тире из разделителей элементов в библиографии:
%\renewcommand*{\newblockpunct}{%
%    \addperiod\space\bibsentence}%block punct.,\bibsentence is for vol,etc.
%%% Изменение точки с запятой на запятую в перечислении библиографических
%%% ссылок:
%\renewcommand*{\multicitedelim}{\addcomma\space}

%%% Возвращаем запись «Режим доступа» %%%
%\DefineBibliographyStrings{english}{%
%    urlfrom = {Mode of access}
%}
%\DeclareFieldFormat{url}{\bibstring{urlfrom}\addcolon\space\url{#1}}

%%% В списке литературы обозначение одной буквой диапазона страниц англоязычного источника %%%
\DefineBibliographyStrings{english}{%
    pages = {p\adddot} %заглавность буквы затем по месту определяется работой самого biblatex
}

%%% В ссылке на источник в основном тексте с указанием конкретной страницы обозначение одной большой буквой %%%
%\DefineBibliographyStrings{russian}{%
%    page = {C\adddot}
%}

%%% Исправление длины тире в диапазонах %%%
% \cyrdash --- тире «русской» длины, \textendash --- en-dash
\DefineBibliographyExtras{russian}{%
  \protected\def\bibrangedash{%
    \cyrdash\penalty\value{abbrvpenalty}}% almost unbreakable dash
  \protected\def\bibdaterangesep{\bibrangedash}%тире для дат
}
\DefineBibliographyExtras{english}{%
  \protected\def\bibrangedash{%
    \cyrdash\penalty\value{abbrvpenalty}}% almost unbreakable dash
  \protected\def\bibdaterangesep{\bibrangedash}%тире для дат
}

%Set higher penalty for breaking in number, dates and pages ranges
\setcounter{abbrvpenalty}{10000} % default is \hyphenpenalty which is 12

%Set higher penalty for breaking in names
\setcounter{highnamepenalty}{10000} % If you prefer the traditional BibTeX behavior (no linebreaks at highnamepenalty breakpoints), set it to ‘infinite’ (10 000 or higher).
\setcounter{lownamepenalty}{10000}

%%% Set low penalties for breaks at uppercase letters and lowercase letters
%\setcounter{biburllcpenalty}{500} %управляет разрывами ссылок после маленьких букв RTFM biburllcpenalty
%\setcounter{biburlucpenalty}{3000} %управляет разрывами ссылок после больших букв, RTFM biburlucpenalty

%%% Список литературы с красной строки (без висячего отступа) %%%
%\defbibenvironment{bibliography} % переопределяем окружение библиографии из gost-numeric.bbx пакета biblatex-gost
%  {\list
%     {\printtext[labelnumberwidth]{%
%       \printfield{prefixnumber}%
%       \printfield{labelnumber}}}
%     {%
%      \setlength{\labelwidth}{\labelnumberwidth}%
%      \setlength{\leftmargin}{0pt}% default is \labelwidth
%      \setlength{\labelsep}{\widthof{\ }}% Управляет длиной отступа после точки % default is \biblabelsep
%      \setlength{\itemsep}{\bibitemsep}% Управление дополнительным вертикальным разрывом между записями. \bibitemsep по умолчанию соответствует \itemsep списков в документе.
%      \setlength{\itemindent}{\bibhang}% Пользуемся тем, что \bibhang по умолчанию принимает значение \parindent (абзацного отступа), который переназначен в styles.tex
%      \addtolength{\itemindent}{\labelwidth}% Сдвигаем правее на величину номера с точкой
%      \addtolength{\itemindent}{\labelsep}% Сдвигаем ещё правее на отступ после точки
%      \setlength{\parsep}{\bibparsep}%
%     }%
%      \renewcommand*{\makelabel}[1]{\hss##1}%
%  }
%  {\endlist}
%  {\item}

%%% Макросы автоматического подсчёта количества авторских публикаций.
% Печатают невидимую (пустую) библиографию, считая количество источников.
% http://tex.stackexchange.com/a/66851/79756
%
\makeatletter
    \newtotcounter{citenum}
    \defbibenvironment{counter}
        {\setcounter{citenum}{0}\renewcommand{\blx@driver}[1]{}} % begin code: убирает весь выводимый текст
        {} % end code
        {\stepcounter{citenum}} % item code: cчитает "печатаемые в библиографию" источники

    \newtotcounter{citeauthorvak}
    \defbibenvironment{countauthorvak}
        {\setcounter{citeauthorvak}{0}\renewcommand{\blx@driver}[1]{}}
        {}
        {\stepcounter{citeauthorvak}}

    \newtotcounter{citeauthorscopus}
    \defbibenvironment{countauthorscopus}
        {\setcounter{citeauthorscopus}{0}\renewcommand{\blx@driver}[1]{}}
        {}
        {\stepcounter{citeauthorscopus}}

    \newtotcounter{citeauthorwos}
    \defbibenvironment{countauthorwos}
        {\setcounter{citeauthorwos}{0}\renewcommand{\blx@driver}[1]{}}
        {}
        {\stepcounter{citeauthorwos}}

    \newtotcounter{citeauthorother}
    \defbibenvironment{countauthorother}
        {\setcounter{citeauthorother}{0}\renewcommand{\blx@driver}[1]{}}
        {}
        {\stepcounter{citeauthorother}}

    \newtotcounter{citeauthorconf}
    \defbibenvironment{countauthorconf}
        {\setcounter{citeauthorconf}{0}\renewcommand{\blx@driver}[1]{}}
        {}
        {\stepcounter{citeauthorconf}}

    \newtotcounter{citeauthor}
    \defbibenvironment{countauthor}
        {\setcounter{citeauthor}{0}\renewcommand{\blx@driver}[1]{}}
        {}
        {\stepcounter{citeauthor}}

    \newtotcounter{citeauthorvakscopuswos}
    \defbibenvironment{countauthorvakscopuswos}
        {\setcounter{citeauthorvakscopuswos}{0}\renewcommand{\blx@driver}[1]{}}
        {}
        {\stepcounter{citeauthorvakscopuswos}}

    \newtotcounter{citeauthorscopuswos}
    \defbibenvironment{countauthorscopuswos}
        {\setcounter{citeauthorscopuswos}{0}\renewcommand{\blx@driver}[1]{}}
        {}
        {\stepcounter{citeauthorscopuswos}}

    \newtotcounter{citeregistered}
    \defbibenvironment{countregistered}
        {\setcounter{citeregistered}{0}\renewcommand{\blx@driver}[1]{}}
        {}
        {\stepcounter{citeregistered}}

    \newtotcounter{citeauthorpatent}
    \defbibenvironment{countauthorpatent}
        {\setcounter{citeauthorpatent}{0}\renewcommand{\blx@driver}[1]{}}
        {}
        {\stepcounter{citeauthorpatent}}

    \newtotcounter{citeauthorprogram}
    \defbibenvironment{countauthorprogram}
        {\setcounter{citeauthorprogram}{0}\renewcommand{\blx@driver}[1]{}}
        {}
        {\stepcounter{citeauthorprogram}}

    \newtotcounter{citeexternal}
    \defbibenvironment{countexternal}
        {\setcounter{citeexternal}{0}\renewcommand{\blx@driver}[1]{}}
        {}
        {\stepcounter{citeexternal}}
\makeatother

\defbibheading{nobibheading}{} % пустой заголовок, для подсчёта публикаций с помощью невидимой библиографии
\defbibheading{pubgroup}{\section*{#1}} % обычный стиль, заголовок-секция
\defbibheading{pubsubgroup}{\noindent\textbf{#1}} % для подразделов "по типу источника"

%%%Сортировка списка литературы Русский-Английский (предварительно удалить dissertation.bbl) (начало)
%%%Источник: https://github.com/odomanov/biblatex-gost/wiki/%D0%9A%D0%B0%D0%BA-%D1%81%D0%B4%D0%B5%D0%BB%D0%B0%D1%82%D1%8C,-%D1%87%D1%82%D0%BE%D0%B1%D1%8B-%D1%80%D1%83%D1%81%D1%81%D0%BA%D0%BE%D1%8F%D0%B7%D1%8B%D1%87%D0%BD%D1%8B%D0%B5-%D0%B8%D1%81%D1%82%D0%BE%D1%87%D0%BD%D0%B8%D0%BA%D0%B8-%D0%BF%D1%80%D0%B5%D0%B4%D1%88%D0%B5%D1%81%D1%82%D0%B2%D0%BE%D0%B2%D0%B0%D0%BB%D0%B8-%D0%BE%D1%81%D1%82%D0%B0%D0%BB%D1%8C%D0%BD%D1%8B%D0%BC
%\DeclareSourcemap{
%    \maps[datatype=bibtex]{
%        \map{
%            \step[fieldset=langid, fieldvalue={tempruorder}]
%        }
%        \map[overwrite]{
%            \step[fieldsource=langid, match=russian, final]
%            \step[fieldsource=presort,
%            match=\regexp{(.+)},
%            replace=\regexp{aa$1}]
%        }
%        \map{
%            \step[fieldsource=langid, match=russian, final]
%            \step[fieldset=presort, fieldvalue={az}]
%        }
%        \map[overwrite]{
%            \step[fieldsource=langid, notmatch=russian, final]
%            \step[fieldsource=presort,
%            match=\regexp{(.+)},
%            replace=\regexp{za$1}]
%        }
%        \map{
%            \step[fieldsource=langid, notmatch=russian, final]
%            \step[fieldset=presort, fieldvalue={zz}]
%        }
%        \map{
%            \step[fieldsource=langid, match={tempruorder}, final]
%            \step[fieldset=langid, null]
%        }
%    }
%}
%Сортировка списка литературы (конец)

%%% Создание команд для вывода списка литературы %%%
\newcommand*{\insertbibliofull}{
    \printbibliography[keyword=bibliofull,section=0,title=\bibtitlefull]
    \ifnumequal{\value{draft}}{0}{
      \printbibliography[heading=nobibheading,env=counter,keyword=bibliofull,section=0]
    }{}
}
\newcommand*{\insertbiblioauthor}{
    \printbibliography[heading=pubgroup, section=0, filter=papersregistered, title=\bibtitleauthor]
}
\newcommand*{\insertbiblioauthorimportant}{
    \printbibliography[heading=pubgroup, section=2, filter=papersregistered, title=\bibtitleauthorimportant]
}

% Вариант вывода печатных работ автора, с группировкой по типу источника.
% Порядок команд `\printbibliography` должен соответствовать порядку в файле common/characteristic.tex
\newcommand*{\insertbiblioauthorgrouped}{
    \section*{\bibtitleauthor}
    \ifsynopsis
    \printbibliography[heading=pubsubgroup, section=0, keyword=biblioauthorvak,    title=\bibtitleauthorvak,resetnumbers=true] % Работы автора из списка ВАК (сброс нумерации)
    \else
    \printbibliography[heading=pubsubgroup, section=0, keyword=biblioauthorvak,    title=\bibtitleauthorvak,resetnumbers=false] % Работы автора из списка ВАК (сквозная нумерация)
    \fi
    \printbibliography[heading=pubsubgroup, section=0, keyword=biblioauthorwos,    title=\bibtitleauthorwos,resetnumbers=false]% Работы автора, индексируемые Web of Science
    \printbibliography[heading=pubsubgroup, section=0, keyword=biblioauthorscopus, title=\bibtitleauthorscopus,resetnumbers=false]% Работы автора, индексируемые Scopus
    \printbibliography[heading=pubsubgroup, section=0, keyword=biblioauthorpatent, title=\bibtitleauthorpatent,resetnumbers=false]% Патенты
    \printbibliography[heading=pubsubgroup, section=0, keyword=biblioauthorprogram,title=\bibtitleauthorprogram,resetnumbers=false]% Программы для ЭВМ
    \printbibliography[heading=pubsubgroup, section=0, keyword=biblioauthorconf,   title=\bibtitleauthorconf,resetnumbers=false]% Тезисы конференций
    \printbibliography[heading=pubsubgroup, section=0, keyword=biblioauthorother,  title=\bibtitleauthorother,resetnumbers=false]% Прочие работы автора
}

\newcommand*{\insertbiblioexternal}{
    \printbibliography[heading=pubgroup,    section=0, keyword=biblioexternal,     title=\bibtitlefull]
}
   % Реализация пакетом biblatex через движок biber
%}

% Вывести информацию о выбранных опциях в лог сборки
\typeout{Selected options:}
\typeout{Draft mode: \arabic{draft}}
\typeout{Font: \arabic{fontfamily}}
\typeout{AltFont: \arabic{usealtfont}}
\typeout{Bibliography backend: \arabic{bibliosel}}
\typeout{Precompile images: \arabic{imgprecompile}}
% Вывести информацию о версиях используемых библиотек в лог сборки
\listfiles

\begin{document}

\thispagestyle{empty}

\noindent%
\begin{tabularx}{\textwidth}{@{}lXr@{}}%
    & & \large{На правах рукописи} \\
    \IfFileExists{images/logo.pdf}{\includegraphics[height=2.5cm]{logo}}{\rule[0pt]{0pt}{2.5cm}}  & &
    \ifnumequal{\value{showperssign}}{0}{%
        \rule[0pt]{0pt}{1cm}
    }{
        \includegraphics[height=1cm]{personal-signature.png}
    }\\
\end{tabularx}

\vspace{0pt plus1fill} %число перед fill = кратность относительно некоторого расстояния fill, кусками которого заполнены пустые места
\begin{center}
    \textbf {\large \thesisAuthor}
\end{center}

\vspace{0pt plus3fill} %число перед fill = кратность относительно некоторого расстояния fill, кусками которого заполнены пустые места
\begin{center}
    \textbf {\Large %\MakeUppercase
    \thesisTitle}

    \vspace{0pt plus3fill} %число перед fill = кратность относительно некоторого расстояния fill, кусками которого заполнены пустые места
    {\large Специальность \thesisSpecialtyNumber\ "---\par <<\thesisSpecialtyTitle>>}

    \ifdefined\thesisSpecialtyTwoNumber
    {\large Специальность \thesisSpecialtyTwoNumber\ "---\par <<\thesisSpecialtyTwoTitle>>}
    \fi

    \vspace{0pt plus1.5fill} %число перед fill = кратность относительно некоторого расстояния fill, кусками которого заполнены пустые места
    \Large{Автореферат}\par
    \large{диссертации на соискание учёной степени\par \thesisDegree}
\end{center}

\vspace{0pt plus4fill} %число перед fill = кратность относительно некоторого расстояния fill, кусками которого заполнены пустые места
{\centering\thesisCity~--- \thesisYear\par}

\newpage
% оборотная сторона обложки
\thispagestyle{empty}
\noindent Работа выполнена в {\thesisInOrganization}.

\vspace{0.008\paperheight plus1fill}
\noindent%
\begin{tabularx}{\textwidth}{@{}lX@{}}
    Научный руководитель: & \textbf{\supervisorFio}\par
    \supervisorRegalia
    \vspace{0.013\paperheight} \\
    Официальные оппоненты: &
    \textbf{\opponentOneFio,}\par
    \opponentOneRegalia,\par
    \opponentOneJobPlace,\par
    \opponentOneJobPost.\par
    \vspace{0.01\paperheight}
    \textbf{\opponentTwoFio,}\par
    \opponentTwoRegalia,\par
    \opponentTwoJobPlace,\par
    \opponentTwoJobPost.\par

    \vspace{0.013\paperheight} \\
    \ifdefined\leadingOrganizationTitle
    Ведущая организация: &
    \ifnumequal{\value{showopplead}}{0}{\vspace{6\onelineskip plus1fill}}{%
        \leadingOrganizationTitle
    }%
    \fi
\end{tabularx}
\vspace{0.008\paperheight plus1fill}

\noindent Защита состоится \defenseDate~на~заседании диссертационного совета \defenseCouncilNumber~при \defenseCouncilTitle~по адресу: \defenseCouncilAddress.

\vspace{0.008\paperheight plus1fill}
\noindent С диссертацией можно ознакомиться в библиотеке \synopsisLibrary.

\vspace{0.008\paperheight plus1fill}
\noindent Отзывы на автореферат в двух экземплярах, заверенные печатью учреждения, просьба направлять по адресу: \defenseCouncilAddress, ученому секретарю диссертационного совета~\defenseCouncilNumber.

\vspace{0.008\paperheight plus1fill}
\noindent{Автореферат разослан \synopsisDate.}

\noindent Телефон для справок: \defenseCouncilPhone.

\vspace{0.008\paperheight plus1fill}
\noindent%
\begin{tabularx}{\textwidth}{@{}%
        >{\raggedright\arraybackslash}b{18em}@{}
        >{\centering\arraybackslash}X
    r
    @{}}
    Ученый секретарь\par
    диссертационного совета\par
    \defenseCouncilNumber,\par
    \defenseSecretaryRegalia
    &
    \ifnumequal{\value{showsecrsign}}{0}{}{%
        \includegraphics[width=2cm]{secretary-signature.png}%
    }%
    &
    \defenseSecretaryFio
\end{tabularx}
        % Титульный лист
%\mainmatter                   % В том числе начинает нумерацию страниц арабскими цифрами с единицы
\mainmatter*                  % Нумерация страниц не изменится, но начнётся с новой страницы
\pdfbookmark{Общая характеристика работы}{characteristic}             % Закладка pdf
\section*{Общая характеристика работы}

\newcommand{\actuality}{\pdfbookmark[1]{Актуальность}{actuality}\underline{\textbf{\actualityTXT}}}
\newcommand{\progress}{\pdfbookmark[1]{Разработанность темы}{progress}\underline{\textbf{\progressTXT}}}
\newcommand{\aim}{\pdfbookmark[1]{Цели}{aim}\underline{{\textbf\aimTXT}}}
\newcommand{\tasks}{\pdfbookmark[1]{Задачи}{tasks}\underline{\textbf{\tasksTXT}}}
\newcommand{\aimtasks}{\pdfbookmark[1]{Цели и задачи}{aimtasks}\aimtasksTXT}
\newcommand{\novelty}{\pdfbookmark[1]{Научная новизна}{novelty}\underline{\textbf{\noveltyTXT}}}
\newcommand{\influence}{\pdfbookmark[1]{Практическая значимость}{influence}\underline{\textbf{\influenceTXT}}}
\newcommand{\methods}{\pdfbookmark[1]{Методология и методы исследования}{methods}\underline{\textbf{\methodsTXT}}}
\newcommand{\defpositions}{\pdfbookmark[1]{Положения, выносимые на защиту}{defpositions}\underline{\textbf{\defpositionsTXT}}}
\newcommand{\reliability}{\pdfbookmark[1]{Достоверность}{reliability}\underline{\textbf{\reliabilityTXT}}}
\newcommand{\probation}{\pdfbookmark[1]{Апробация}{probation}\underline{\textbf{\probationTXT}}}
\newcommand{\contribution}{\pdfbookmark[1]{Личный вклад}{contribution}\underline{\textbf{\contributionTXT}}}
\newcommand{\publications}{\pdfbookmark[1]{Публикации}{publications}\underline{\textbf{\publicationsTXT}}}

%    {\actuality}
    Под сложным теплообменом понимают
    процесс распространения тепла излучением, теплопроводностью и
    конвекцией.
    При этом важно, что радиационный перенос тепла дает
    существенный вклад в распределение температурного поля.
    Для многих инженерных приложений моделирование процессов теплопередачи имеет
    первостепенное значение (например, это относится к радиационному
    теплообмену в полупрозрачных материалах).
    Температура – важнейший параметр практически во всех производственных этапах.
    Следовательно, знание температурных полей необходимо для контроля производственных
    процессов и качества результата.

    С математической точки зрения процесс сложного теплообмена
    моделируется системой уравнений с частными производными, включающей
    уравнение теплопроводности, а также интегро-дифференциальное уравнение
    переноса теплового излучения~\cite{Ozisik1976, Sparrow1971, howell2010thermal, modest2013radiative}.
    Теоретический и численный анализ различных задач для такой системы является весьма
    затратным поскольку интенсивность излучения зависит не только от
    пространственной и временной переменных, но и от направления
    распространения излучения.
    Имеется ряд аппроксимаций уравнения переноса излучения, в том числе диффузионное
    $P_1$-приближение, в котором интенсивность излучения усредняется по направлениям.
    $P_1$-приближение является частным случаем метода сферических гармоник ($PN$-приближения)
    и упрощенного метода сферических гармоник ($SPN$-приближения).
    В данной работе рассмотрены модели, основанные на $P_1$-приближении.


    В классических прямых задачах сложного теплообмена задаются
    параметры системы, и по ним вычисляется состояние системы –
    температурное поле и интенсивность теплового излучения.
    Обратные задачи сложного теплообмена состоят в разыскании неизвестных параметров
    системы по некоторой дополнительной информации (условия переопределения)
    о температурном поле или интенсивности излучения.
    Методы, основанные на решении обратных задач, продолжают оставаться
    актуальным, быстро развивающимся направлением исследований.
    Отметим трудности, возникающие при решении обратных задач сложного
    теплообмена.


    Эти задачи математически классифицируются как некорректные в общем
    смысле из-за неустойчивости решений.
    Развитие методов оптимизации позволяет исправить проблемы некорректности исследуемых задач.
    В основе таких методов лежит идея замены исходной задачи на задачу оптимизации и
    фактически использует методику А.Н. Тихонова регуляризации задачи.

    Диссертация посвящена теоретическому и численному анализу одному из
    наиболее распространенных на практике типу обратных задач – граничных
    обратных задач сложного теплообмена в трёхмерной области в рамках
    $P_1$-приближения уравнения переноса излучения.
    При этом основное внимание уделяется разработке
    и обоснованию оптимизационных методов решения указанных задач.
    Исследование обратных задач для математических моделей
    радиационного теплопереноса, учитывающих одновременно вклад эффектов
    теплопроводности и излучения, даёт теоретическую основу для эффективных
    инженерных решений в различных областях, таких как производство стекла,
    лазерная термотерапия и другие.


    \textbf{Степень разработанности темы исследования.}
    Имеется значительное число работ, посвященных теоретическому анализу моделей сложного
    теплообмена.
    Результаты А.\ А.\ Амосова, M.\ Laitinen, T.\ Tiihonen, P.-E.\ Druet и др.~\cite{
        Amosov2005,
        Amosov2009,
        Amosov2010nonstationary,
        amosov2010stationary,
        druet2009weak,
        druet2010weak,
        laitinen2001conductive,
        metzger1999existence,
        Philip2010,
        Tiihonen1997a,
        Tiihonen1997b,
        amosov17,
        amosov18,
        Amosov20,
        Amosov21-1,
        Amosov2021-2,
        Amosov23} %[Гренкин 1, 2, 30, 31, 56, 57, 80, 86, 95, 119, 120 Добавить свежие статьи Амосова]
    посвящены анализу разрешимости моделей сложного теплообмена,
    которые включают уравнение теплопроводности с нелинейным нелокальным
    краевым условием, моделирующим тепловое излучение границы области и
    теплообмен излучением между частями границы,
    в~\cite{Amosov2009, amosov2010stationary, druet2009weak, laitinen2001conductive,
        Philip2010, Tiihonen1997a, Tiihonen1997b} %[2, 31, 56, 80, 95, 119, 120]
    рассмотрены стационарные модели,
    в~\cite{
        Amosov2005, Amosov2010nonstationary,
        druet2010weak, laitinen2001conductive,
        metzger1999existence,
        Philip2010,
        Tiihonen1997a} — нестационарные.%[1, 30, 57, 80, 86, 95, 119]

    Анализ моделей сложного теплообмена в полупрозрачной среде, в
    которых для моделирования радиационного теплообмена используется
    полное уравнение переноса излучения представлен в следующих статьях.
    В~\cite{asllanaj2003existence, kelley1996existence} доказана
    однозначная разрешимость одномерных стационарных
    задач радиационно-кондуктивного теплообмена,
    в~\cite{ghattassi2018existence, Porzio2004, Thompson2004} доказана
    однозначная разрешимость трехмерных задач: в~\cite{Thompson2004} исследована
    стационарная модель, в~\cite{Porzio2004} -- нестационарная,
    в~\cite{ghattassi2018existence} -- квазистационарная.
    Под квазистационарными моделями сложного теплообмена понимаются
    модели, включающие нестационарное уравнение теплопроводности и
    стационарное уравнение переноса излучения.
    В работах R.\ Pinnau, O.\ Tse~\cite{Pinnau2007, Pinnau2013}
    проведен теоретический анализ квазистационарных моделей сложного
    теплообмена на основе $SP_1$- и $SP_3$-приближений.
    Эти модели включают уравнение теплопроводности, стационарное $SP_N$-приближение, а также
    в~\cite{Pinnau2013} уравнения Навье–Стокса в приближении Буссинеска.
    В работах Г.\ В.\ Гренкина, А.\ Е.\ Ковтанюка,
    А.\ Ю.\ Чеботарева~\cite{Kovtanyuk2014, Kovtanyuk2016, Chebotarev22, Pak23}
    доказана однозначная разрешимость краевых задач для моделей
    сложного теплообмена на основе $P_1$-приближения, доказана сходимость
    метода простой итерации нахождения решения.
    В~\cite{Chebotarev2016Odnaznachnaya} доказана однозначная
    разрешимость субдифференциальной краевой задачи с многозначной
    зависимостью коэффициента излучения границы от интенсивности
    излучения.

    Подробный обзор различных численных методов решения задач сложного
    теплообмена можно найти в~\cite{modest2013radiative}.
    Большинство методов можно разделить на четыре класса:
    \begin{itemize}
        \item Методы диффузионной аппроксимации.
        Типичным примером является приближение Розеланда~\cite{farina2011mathematical, Siedow2011}.
        \item Методы, в которых зависимость теплового излучения от направления
        выражается рядом специальных базисных функций.
        Представителями методов являются $PN$ и $SPN$-аппроксимации.
        \item Методы, использующие дискретизацию угловой зависимости, например
        метод дискретных ординат, который является наиболее популярным.
        \item Методы Монте-Карло.
    \end{itemize}


    Статьи~\cite{kovtanyuk2013iterative, Thommes2002, Pinnau2008, Siewert1991}
    посвящены численному моделированию в рамках
    диффузионных моделей сложного теплообмена, в~\cite{gallouet2016analysis} исследована схема
    метода конечных объемов для решения квазистационарной системы
    уравнений сложного теплообмена на основе
    $P_1$-приближения уравнения переноса излучения.
    Сравнение $P_1$-приближения с другими методами
    аппроксимации уравнения переноса излучения проводилось
    в~\cite{modest2014elliptic, frank2011adaptive, kovtanyuk2012, Thommes2002, Larsen2002, Frank2007}.
    Численный анализ нестационарного $P_1$-приближения выполнен
    в~\cite{Addam2015, olbrant2013asymptotic, frank2010optimal, frank2011adaptive, Frank2007}.

    Основополагающие работы А.Н. Тихонова~\cite{TikhonovSamarskii1972} и его коллег, которые
    исследовали обратные задачи для уравнений математической физики, стали
    отправной точкой для разработки методов преодоления неустойчивостей в обратных задачах.
    О.\ М.\ Алифанов~\cite{Aliphanov77, Aliphanov99, Aliphanov2009}
    в своих работах, посвященных обратным задачам сложного теплообмена, также
    сделал значительный вклад в развитие этой области.
    J.\ V.\ Beck~\cite{Beck1985-fg} предложил
    новые подходы к решению обратных задач, основанных на методах
    оптимизации и статистическом анализе.
    Отметим также серьезный теоретический анализ обратных задач тепломассопереноса
    представленный в работах С.\ Г.\ Пяткова и его учеников~\cite{Pyatkov19, Pyatkov22, Pyatkov23}.


    Решение обратных задач оптимизационным методом можно
    рассматривать как соответствующую задачу оптимального управления
    системой, моделируемой уравнениями сложного теплообмена.
    Обратно, каждую задачу оптимального управления можно интерпретировать как
    обратную экстремальную задачу.
    Отметим следующие работы, посвященные
    решению обратных задач в рамках моделей сложного теплообмена с полным
    уравнением переноса излучения.
    В~\cite{end2011analytical} проведен теоретический анализ
    задачи оптимального управления источниками тепла в рамках
    квазистационарной модели сложного теплообмена, включающей полное
    уравнение переноса излучения.
    В~\cite{end2010optimization} разработан численный алгоритм
    решения задачи оптимального управления источниками тепла и излучения в
    рамках стационарной модели сложного теплообмена с полным уравнением
    переноса излучения.
    Работа~\cite{Pereverzyev2008} посвящена теоретическому и численному
    анализу обратной задачи восстановления начального распределения
    температуры по известной зависимости температуры на границе области от
    времени в рамках квазистационарной модели сложного теплообмена.
    Отметим также работы~\cite{birgelis2003optimal, meyer2006optimal, meyer2009state, Philip2010},
%    [48, 87, 88, 95],
    посвященные анализу задач оптимального управления для стационарных
    моделей сложного теплообмена в прозрачной среде, включающих уравнение теплопроводности с
    нелинейным нелокальным краевым условием, моделирующим тепловое
    излучение границы области и теплообмен излучением между частями
    границы, и работу~\cite{belmiloudi2014nonlinear},
    в которой построен численный алгоритм решения
    задачи оптимального управления граничными коэффициентами
    в одномерной нестационарной модели, включающей уравнение
    теплопроводности с нелинейным краевым условием, которое описывает
    тепловое излучение границ.
    В~\cite{clever2012optimal, clever2014model, frank2010optimal, lang2005adaptive, Pinnau2007b, Pinnau2004}
%    [50, 51, 64, 81, 98, 100]
    разработаны численные
    методы решения задач оптимального управления граничной температурой
    для квазистационарной модели сложного теплообмена на основе $P_1$-приближения,
    при этом в~\cite{clever2012optimal, clever2014model, lang2005adaptive} использовалась
    модель с учетом зависимости коэффициента поглощения от частоты излучения:
    в~\cite{frank2010optimal, Pinnau2007b}
    минимизировалось отклонение поля температуры от заданного и для
    решения задачи оптимизации применялся метод Ньютона, в~\cite{lang2005adaptive} для решения
    задачи минимизации отклонения поля температуры от заданного применялся
    метод проекции градиента, в~\cite{clever2012optimal, Pinnau2004} минимизировалась норма градиента
    температуры и для решения задачи оптимизации применялся метод проекции
    градиента, в~\cite{clever2014model} решалась задача минимизации отклонения поля
    температуры от заданного на основе серии из трех моделей,
    аппроксимирующих уравнение переноса излучения с разной точностью,
    использовался оптимизационный метод второго порядка.
    В~\cite{Pinnau2004} проведен
    теоретический анализ задачи оптимального управления температурой на
    границе области в рамках стационарной диффузионной модели сложного
    теплообмена, для численного решения задачи управления применен метод
    проекции градиента.
    В работах А.\ Е.\ Ковтанюка, А.\ Ю.\ Чеботарева и др.\ \cite{Kovtanyuk2014,
        Mesenev2018, astrakhantseva2017design, Chebotarev2015,
        Kovtanyuk2014TheoreticalAnalysis, Chebotarev23nonhomogeneus, Chebotarev23optimal}.
    исследованы задачи оптимального
    управления коэффициентом излучения границы области в рамках
    стационарной модели сложного теплообмена на основе $P_1$-приближения.


    В последние годы исследования моделей сложного теплообмена стали
    особенно актуальными в связи с практическими приложениями.
    Ниже приведены некоторые примеры исследований практической применимости
    рассматриваемых моделей.
    В работах~\cite{clever2012optimal, clever2014model,
        frank2010optimal, klar2005, lang2005adaptive, farina2011mathematical,
        Thommes2002, Pinnau2007b, Pinnau2004, Seaid2007, Larsen2002}
    модели сложного теплообмена на основе -приближений
    применялись для моделирования и оптимизации сложного теплообмена при
    производстве стекла, в~\cite{frank2004comparison, seaid2004efficient, Seaid2005}
    моделировался сложный теплообмен в
    камерах сгорания газовых турбин, в~\cite{Backofen2004} диффузионные модели
    использовались для моделирования переноса теплового излучения в
    растущем кристалле.
    Также $P_1$-приближение применялось в составе моделей
    лазерной термотерапии~\cite{Dombrovskii2015, Tse2012, Hubner2017}.
    Процедура внутривенной лазерной
    абляции (ВВЛА) является безопасной и эффективной в лечении варикозных
    вен~\cite{Endovenous_vandenBos2009}.
    Математическое моделирование лучевых и тепловых процессов
    при ВВЛА имеет решающее значение для определения оптимальных
    параметров излучения, обеспечивающих достаточно высокие температуры
    внутри вены для успешной облитерации, обеспечивая при этом сохранность
    окружающих тканей.
    Результаты численного моделирования для различных
    длин волн и диаметров вен обсуждались в ряде
    работ~\cite{van2014optical, Some_Poluektova2014, Endovenous_Malskat2014, Mathematical_Mordon2006}.
    Задачи оптимального управления для модели ВВЛА рассматриваются
    в~\cite{Optimal_Kovtanyuk2020, Inverse_Kovtanyuk2021}.
    В~\cite{Optimal_Kovtanyuk2020} поставлена задача оптимального
    управления для модели реакция-диффузия,
    описывающей процедуру ВВЛА, которая заключается в
    аппроксимации заданного температурного профиля в определенной точке
    области модели.
    В~\cite{Inverse_Kovtanyuk2021} изучается аналогичная задача оптимального управления.
    Здесь целевой функционал выбран таким образом, что его
    минимизация позволяет достичь заданного распределения температуры в
    различных частях области модели.


    Несмотря на представленный в обзоре значительный объем исследований,
    включающий анализ обратных и обратных экстремальных задач, ряд важных
    задач, связанных с анализом корректности стационарных,
    квазистационарных и квазилинейных \textit{\textbf{моделей сложного теплообмена}},
    корректности постановок поиска квазирешений граничных обратных задач,
    \textit{\textbf{построением и обоснованием сходимости}} оптимизационных алгоритмов
    решения обратных задач и задач с краевыми условиями Коши, а также с
    \textit{\textbf{разработкой комплекса программ для проведения вычислительных экспериментов}}
    и тестирования предложенных алгоритмов оставался нерешенным.
    Настоящая работа посвящена решению указанных проблем.


    \textbf{Цели и задачи диссертационной работы.}
    \textit{Целью работы} является является теоретический и численный анализ граничных
    обратных задач, включая задачи с условиями типа Коши на границе области,
    и задач оптимального управления для моделей сложного теплообмена на
    основе $P_1$-приближения уравнения переноса излучения, а именно:
    \begin{itemize}[leftmargin=5.5mm]
        \renewcommand\labelitemi{--}
        \item исследование разрешимости краевых и начально-краевых задач для
        квазистационарных и квазилинейных моделей сложного теплообмена;
        \item разработка оптимизационных алгоритмов решения обратных задач и задач
        с краевыми условиями Коши, теоретический анализ и обоснование их
        сходимости;
        \item разработка комплекса программ для проведения вычислительных
        экспериментов и тестирования предложенных алгоритмов.
    \end{itemize}
    \textit{Для достижения целей работы были сформулированы следующие задачи:}
    \begin{itemize}[leftmargin=5.5mm]
        \renewcommand\labelitemi{--}
        \item доказать существование и единственность решения начально-краевой
        задачи для квазистационарной и квазилинейной моделей сложного теплообмена,
        разработать итерационный алгоритм нахождения решения и обосновать его
        сходимость;
        \item получить условия существования квазирешения обратной задачи с
        неизвестным коэффициентом отражения на границе, вывести условия
        оптимальности первого порядка и получить численный алгоритм;
        \item выполнить анализ оптимизационных методов решения задач сложного
        теплообмена с условиями Коши на границе для стационарной и
        квазистационарной моделей, исследовать разрешимость возникающих
        регуляризованных экстремальных задач, получить условия оптимальности,
        обосновать сходимость решений регуляризованных задач к решению задач с
        условиями Коши на границе при стремлении параметра регуляризации к
        нулю;
        \item провести теоретический анализ обратных экстремальных задач с фазовыми
        ограничениями для квазилинейной модели сложного теплообмена,
        рассмотреть аппроксимации задачами с штрафными функционалами и
        доказать сходимость их решений при увеличении штрафа;
        \item предложить численные методы решения рассматриваемых краевых и
        оптимизационных задач сложного теплообмена, адаптировать метод
        Ньютона и градиентные методы для их решения, реализовать полученные
        методы в виде программных систем, осуществить тестирование
        предложенных алгоритмов используя данные реальных сред и материалов.
    \end{itemize}


    \textbf{Научная новизна.}
    В работе получены новые априорные оценки решений
    начально-краевых задач для квазистационарных и квазилинейных уравнений
    сложного теплообмена и доказана их нелокальная однозначная
    разрешимость.
    Для рассмотренных моделей сложного теплообмена
    рассмотрены новые постановки граничных обратных задач, предложены
    оптимизационные методы их решения.
    Выполнен теоретический анализ возникающих новых экстремальных задач.
    Представлены априорные оценки решений регуляризованных задач и впервые
    обоснована сходимость их решений к точным решениям обратных задач.
    Для решения задач с фазовыми ограничениями, предложены алгоритмы,
    основанные на аппроксимации экстремальными задачами со штрафом.
    Разработаны и протестированы новые алгоритмы решения прямых,
    обратных и экстремальных задач для моделей сложного теплообмена.


    \textbf{Теоретическая и практическая значимость.}
    Научная значимость результатов диссертации основана, с одной
    стороны, на решении открытых задач, связанных с корректностью моделей
    теплообмена, учитывающих тепловое излучение, теоретическим и
    численным анализом граничных обратных задач и задач оптимального
    управления.
    С другой стороны, развитие новых оптимизационных методов
    решения рассматриваемых нелинейных обратных задач является основой для
    анализа прикладных задач и задач проектирования систем с заданными
    экстремальными свойствами.
    Решение граничных обратных задач имеет практическое значение при
    выборе оптимальных параметров границы области, в которой происходит
    процесс сложного теплообмена.
    Задачи оптимизации имеют важное
    практическое применение при выборе параметров системы для получения
    необходимого распределения температуры или уровня теплового излучения.
    Необходимость выбора параметров системы возникает при проектировании
    инженерных установок, в которых присутствуют процессы сложного
    теплообмена.
    Теоретическая значимость работы обусловлена также тем, что
    результаты, связанные с корректностью рассматриваемых задач,
    сходимостью предлагаемых численных алгоритмов имеют нелокальный
    характер и не содержат нефизичных ограничений типа малости
    определенных параметров.
    Представленный комплекс программ имеет открытый характер и может
    дополняться для решения прикладных обратных задач сложного
    теплообмена.
    С помощью разработанного комплекса программ проведено численное
    моделирование процесса внутривенной лазерной абляции, изучено влияние
    различных факторов на распределения полей температуры и интенсивности
    излучения.


    \textbf{Методология и методы исследования.}
    В работе широко использовались методы математического и функционального анализа,
    теории дифференциальных уравнений в частных производных, теории экстремальных задач.
    Для разработки численных алгоритмов решения применялись методы
    вычислительной математики, объектно-ориентированное и функциональное
    программирование, методы оптимизации и другие.
    Методология исследования обратных задач, рассмотренных в работе,
    заключается в следующем.
    Предварительно выполняется теоретический анализ краевых или начально-краевых задач, моделирующих
    рассматриваемый процесс.
    Для построения оптимизационного алгоритма ставится экстремальная задача или задача
    оптимального управления для рассматриваемой модели сложного теплообмена и проводится ее
    теоретический анализ, включающий нахождение условий разрешимости и
    вывод условий оптимальности.
    Для регуляризованных задач обосновывается
    сходимость последовательности их решений.
    Далее, на основе полученных условий оптимальности
    строится градиентный метод решения обратной задачи.
    Полученные алгоритмы реализуются в виде комплекса программ и
    теоретические результаты тестируются на численных примерах,
    иллюстрирующих эффективность предложенных методов.


    \textbf{Положения, выносимые на защиту.}

    \textit{В области математического моделирования:}

    \newcounter{nameOfYourChoice}
    \begin{enumerate}[leftmargin=5.5mm]
        \item Доказательство однозначной разрешимости начально-краевой задачи,
        моделирующей квазистационарный сложный теплообмен в трехмерной
        области.
        \item Доказательство однозначной разрешимости квазилинейной начально-
        краевой задачи, моделирующей сложный теплообмен с нелинейной
        зависимостью коэффициента теплопроводности от температуры.
        \item Обоснование существования квазирешения обратной задачи с неизвестным
        коэффициентом отражения на части границы и условием переопределения на
        другой части.\ Вывод необходимых условий оптимальности.
        \item Вывод условий разрешимости экстремальных задач, аппроксимирующих
        решения граничных обратных задач (задач с условиями Коши на границе
        области) для стационарной и квазистационарной моделей сложного
        теплообмена.
        \item Построение систем оптимальности и доказательство их невырожденности
        для задач оптимального управления стационарными, квазистационарными и
        квазилинейными уравнениями сложного теплообмена.
        \setcounter{nameOfYourChoice}{\value{enumi}}
    \end{enumerate}

    \textit{В области численных методов:}

    \begin{enumerate}[leftmargin=5.5mm]
        \setcounter{enumi}{\value{nameOfYourChoice}}
        \item Разработка численного алгоритма решения квазилинейной начально-
        краевой задачи, моделирующей сложный теплообмен и доказательство его
        сходимости.
        \item Обоснование сходимости последовательности решений задач
        оптимального управления к решениям задач с условиями Коши на границе
        при стремлении параметра регуляризации к нулю.
        \item Обоснование сходимости алгоритма решения экстремальных обратных
        задач с ограничениями температурных полей методом штрафа.
        \setcounter{nameOfYourChoice}{\value{enumi}}
    \end{enumerate}

    \textit{В области комплексов программ:}

    \begin{enumerate}[leftmargin=5.5mm]
        \setcounter{enumi}{\value{nameOfYourChoice}}
        \item Разработка программ, реализующих численное моделирование процессов
        сложного теплообмена на основе метода конечных элементов.\ Реализация и
        тестирование оптимизационных алгоритмов решения граничных обратных
        задач для стационарных, квазистационарных и квазилинейных моделей.
    \end{enumerate}


    \textbf{Степень достоверности и апробация результатов.}
    Достоверность полученных в диссертации теоретических результатов основывается на
    использовании методов функционального анализа, дифференциальных
    уравнений, теории оптимального управления распределенными системами
    задач и строгих математических доказательствах.
    Достоверность результатов численного моделирования обеспечивается
    доказательством сходимости предложенных итерационных процессов
    и тестированием разработанного комплекса программ.

    Основные результаты диссертации докладывались и обсуждались на
    научных семинарах департамента математического и компьютерного моделирования
    ДВФУ, института прикладной математики ДВО РАН, института автоматики
    и процессов управления ДВО РАН и на следующих научных конференциях:
    \begin{itemize}[leftmargin=5.5mm]
        \item Региональная научно-практическая конференция студентов, аспирантов
        и молодых учёных по естественным наукам (Владивосток, 2018, 2019);
        \item Workshop on Computing Technologies and Applied Mathematics
        (Вычислительные технологии и прикладная математика, Владивосток, 2022);
        \item International Conference DAYS on DIFFRACTION (Санкт-Петербург, 2021, 2023);
        \item International Workshop on Mathematical Modeling and Scientific Computing (Мюнхен, 2020, 2022).
    \end{itemize}

    {\publications}
    Результаты диссертации опубликованы в 10 статьях, из них 4
    статьи~\cite{mesenev_23_opt, mesenev_22_penalty, mesenev_20_alg, mesenev_18_boundary}
    в «Дальневосточном математическом журнале», индексируемом
    в ядре РИНЦ и в международных базах научного цитирования (MathSciNet,
    zbMATH), 1 статья~\cite{mesenev_20_opt_proc} в трудах CEUR (Scopus),
    1 статья~\cite{mesenev_23_math}
    в «Journal of Physics: Conference Series» (Scopus),
    2 статьи~\cite{mesenev_21_optimal_proc, mesenev_23_inv_proc}
    в Proceedings of the International Conference Days on Diffraction (Web of Science, Scopus),
    2 статьи~\cite{mesenev_23_problem, Mesenev_22_analysis}
    в «Журнале вычислительной математики и математической физики»
    (Web of Science, Scopus, MathSciNet, zbMATH).

    Получено 3 свидетельства о регистрации программ для ЭВМ~\cite{progbib1, progbib2, progbib3}.

    \textbf{Личный вклад автора.}
    Результаты в области математического моделирования получены совместно с научным руководителем.
    Результаты в области численных методов и комплексов программ получены автором самостоятельно.
 % Характеристика работы по структуре во введении и в автореферате не отличается (ГОСТ Р 7.0.11, пункты 5.3.1 и 9.2.1), потому её загружаем из одного и того же внешнего файла, предварительно задав форму выделения некоторым параметрам

%Диссертационная работа была выполнена при поддержке грантов \dots

%\underline{\textbf{Объем и структура работы.}} Диссертация состоит из~введения,
%четырех глав, заключения и~приложения. Полный объем диссертации
%\textbf{ХХХ}~страниц текста с~\textbf{ХХ}~рисунками и~5~таблицами. Список
%литературы содержит \textbf{ХХX}~наименование.

\pdfbookmark{Содержание работы}{description}                          % Закладка pdf
\section*{Содержание работы}
Во \underline{\textbf{введении}} обоснована актуальность диссертационной работы,
сформулирована цель и аргументирована научная новизна исследований, показана
теоретическая и практическая значимость полученных результатов,
представлены выносимые на защиту научные положения.


\underline{\textbf{Первая глава}} посвящена анализу
диффузионных моделей сложного теплообмена.
В разделе 1.1 приведена модель сложного теплообмена с
полным уравнением переноса излучения, а также её нормализованный вариант.
Раздел 1.2 посвящен выводу модели переноса излучения в
$P_1$ приближении.
Для получения которого интенсивность излучения и фазовая функция заменяются линейным приближением
соответствующих функций.

\begin{gather*}
    I^{*}(x, \omega, t) = \varphi(x, t)
    +\boldsymbol{\Phi}(x, t) \cdot \omega, \label{eq:1_2:14}\\
    P\left(\omega, \omega^{\prime}\right)= 1
    + A \omega \cdot \omega^{\prime}. \label{eq:1_2:15}
\end{gather*}

Таким образом, диффузионное $P_1$ приближение уравнения переноса излучения
представляется в виде
\begin{equation}
    \label{eq:1_2:23}
    \frac{\partial \theta(x, t)}{\partial t}-a \Delta \theta(x, t)
    +b \kappa_{a}\left(\theta^{4}(x, t)-\varphi(x, t)\right)=0.
\end{equation}
\begin{equation}
    \label{eq:1_2:21}
    - \alpha \Delta \varphi(x, t)+\kappa_{a}\left(\varphi(x, t)-\theta^{4}(x, t)\right)=0.
\end{equation}
\begin{equation}
    \label{eq:1_2:24}
    \alpha \frac{\partial \varphi(x, t)}{\partial n}+
    \gamma(x)\left(\varphi(x, t)-\theta_{b}^{4}(x, t)\right)=0,
\end{equation}
дополняется граничным условием для температуры
\begin{equation}
    \label{eq:1_2:25}
    a \frac{\partial \theta(x, t)}{\partial n}
    +\beta(x)\left(\theta(x, t)-\theta_{b}(x, t)\right)=0,
\end{equation}
а также начальными условиями
\begin{equation}
    \label{eq:1_2:26}
    \theta(x, 0)=\theta_{0}(x), \quad \varphi(x, 0)=\varphi_{0}(x).
\end{equation}

В разделе 1.3 приведены определения функциональных пространств,
вспомогательных утверждений и определений,
которые используются в дальнейшем.

В разделе 1.4 приведены теоретические результаты для
стационарной модели сложного теплообмена
в ограниченной области $\Omega \subset \mathbb{R}^3$ с границей $\Gamma=\partial \Omega$.
\begin{equation}
    \label{eq:1_4:4-2}
    - \alpha \Delta \varphi + \kappa_a \varphi = \kappa_a \theta^4,
\end{equation}
\begin{equation}
    \label{eq:1_4:4-1}
    -a \Delta \theta + \textbf{v} \cdot \nabla \theta
    + b \kappa_a \theta^4 =  b \kappa_a \varphi,
\end{equation}
с граничными условиями
\begin{equation}
    \label{eq:1_4:4-4}
    \alpha \frac{\partial \varphi}{\partial \mathbf{n}} + \gamma
    (\varphi-\theta_b^4)|_{\Gamma} = 0.
\end{equation}
\begin{equation}
    \label{eq:1_4:4-3}
    a \frac{\partial \theta}{\partial \mathbf{n}}
    +\left.\beta\left(\theta-\theta_{b}\right)\right|_{\Gamma}=0,
\end{equation}


Здесь $\theta$ -- нормализованная температура, $\varphi$ --
нормализованная интенсивность излучения, усредненная по всем
направлениям, $\textbf{v}$ -- заданное поле скоростей, $\kappa_a$ --
коэффициент поглощения.
Постоянные $a$, $b$ и $\alpha$
определяются следующим образом:
\[
    a=\frac{k}{\rho c_v},\quad b = \frac{4\sigma n^2 T_{\max}^3}{\rho c_v},
    \quad \alpha=\frac{1}{3\kappa - A \kappa_s},
\]
где $k$ -- теплопроводность, $c_v$ -- удельная теплоемкость, $\rho$ --
плотность, $\sigma$ -- постоянная Стефана-Больцмана, $n$ --
показатель преломления, $T_{\max}$ -- максимальная температура в
ненормализованной модели, $\kappa = \kappa_s + \kappa_a$ -- коэффициент
полного взаимодействия, $\kappa_s$ -- коэффициент рассеяния.

%TODO: Проверить эти формулы!
Раздел 1.5 посвящен квазистационарной модели сложного теплообмена.
Представлена следующая начально-краевая задача:
\begin{equation}
    \label{eq:1_5:1}
    \begin{split}
        & \frac{\partial \theta}{\partial t} - a \Delta \theta
        + b \kappa_{a} \left(|\theta| \theta^{3}-\varphi\right) = 0,\\
        & - \alpha \Delta \varphi
        + \kappa_{a} \left(\varphi-|\theta| \theta^{3}\right) = 0,
        \quad x \in \Omega, \quad 0 < t < T;
    \end{split}
\end{equation}
\begin{align}
    a \left(\partial_{n} \theta+\theta\right)=r,
    & \quad \alpha\left(\partial_{n} \varphi
    + \varphi\right) = u \text { на } \Gamma;  \label{eq:1_5:2}\\
    & \left.\theta\right|_{t=0} = \theta_{0}. \label{eq:1_5:3}
\end{align}
Данная модель описывает систему связанных уравнений в частных про
изводных, моделирующих квазистационарный радиационный и диффузионный
теплообмен.
Дано определение слабого решения, доказана лемма о существовании и единственности
слабого решения, а также справедливости утверждений
\[
    \psi=[\theta]^{5 / 2} \in L^{\infty}(0, T ; H) \cap L^{2}(0, T ; V),
    \quad[\theta]^{4} \in L^{2}(0, T ; H).
\]
Указанный результат используется в параграфе~2.3 в анализе
оптимизационного метода для квазистационарной модели.

В разделе 1.6 рассмотрена квазилинейная модель сложного теплообмена,
представленная начально-краевой задачей в ограниченной трехмерной
области $\Omega$ с отражающей границей $\Gamma=\partial \Omega$:

\begin{equation}
    \label{eq:1_6:1}
    \sigma \partial \theta / \partial t
    -\operatorname{div}(k(\theta) \nabla \theta)
    +b\left(\theta^{3}|\theta|-\varphi\right)=f,
\end{equation}
\begin{equation}
    \label{eq:1_6:2}
    -\operatorname{div}(\alpha \nabla \varphi)
    +\beta\left(\varphi-\theta^{3}|\theta|\right)=g, x \in \Omega, 0<t<T,
\end{equation}
\begin{equation}
    \label{eq:1_6:3}
    k(\theta) \partial_{n} \theta+\left.p\left(\theta-\theta_{b}\right)\right|_{\Gamma}=0,
    \alpha \partial_{n} \varphi
    +\left.\gamma\left(\varphi-\theta_{b}^{4}\right)\right|_{\Gamma}=0,
    \left.\quad \theta\right|_{t=0}=\theta_{in}.
\end{equation}
Приведено определение слабой формулировки задачи~\eqref{eq:1_6:1}--\eqref{eq:1_6:3},
даны априорные оценки решений и доказано существование решения начально-краевой задачи.
Представлены результаты единственности ограниченных решений,
а также сходимость к решению начально-краевой задачи
последовательности, построенной рекурсивно.


\underline{\textbf{Вторая глава}} посвящена исследованию граничных обратных задач для
стационарной и квазистационарной моделей сложного теплообмена.
При решении таких задач часто встречаются краевые условия для
эллиптических или параболических уравнений,
когда на границе (части границы) задаётся неизвестная функция
и её нормальная производная (условия Коши).


В разделе 2.1 рассмотрена нормализованная стационарная модель,
описывающая процесс радиационного теплопереноса в
области $\Omega \subset \mathbb{R}^3$ с липшицевой границей $\Gamma$.
Модель имеет следующий вид
\begin{equation}
    \label{eq:2_1:initial}
    \begin{aligned}
        - a \Delta \theta + b \kappa_a(\theta ^ 3 | \theta | - \varphi) = 0,  \\
        - \alpha \Delta \varphi + \kappa_a (\varphi - \theta ^3 | \theta |) = 0,
    \end{aligned}
\end{equation}
и дополняется граничными условиями на
$\Gamma \coloneqq \partial \Omega =\overline{\Gamma}_0 \cup \overline{\Gamma}_1 \cup \overline{\Gamma}_2$,
где части границы $\Gamma_0, \Gamma_1, \Gamma_2$ не имеют пересечений.

\begin{equation}
    \label{eq:2_1:initial-boundary}
    \begin{aligned}
        \Gamma &: \; a \partial_n \theta + \beta (\theta - \theta _b) = 0, \\
        \Gamma_0 \cup \Gamma_2 &: \; \alpha \partial_n \varphi
        + \gamma(\varphi - \theta_b ^4 ) = 0, \\
        \Gamma_1 &: \; \alpha \partial_n \varphi + u(\varphi - \theta_b ^4 ) = 0. \\
    \end{aligned}
\end{equation}
Функции $\gamma, \theta_b, \beta$ -- являются известными.
Функция $u$ характеризует отражающие свойства участка границы $\Gamma_1$.
Предполагается, что
\begin{equation}
    \label{eq:2_1:control_bounds}
    0 < u_1 \leq u \leq u_2,
\end{equation}
где $u_1$ и $u_2$ -- заданные ограниченные функции.

Рассматриваемая экстремальная задача состоит в минимизации функционала
\begin{equation}
    \label{eq:2_1:quality}
    J(\theta) = \frac{1}{2} \int_{\Gamma_2} (\theta - \theta_0)^2 d\Gamma
\end{equation}
на решениях краевой задачи.

Далее ставится задача нахождения квазирешения, которая состоит в минимизации
функционала~\eqref{eq:2_1:quality} на компоненте $\theta$ решения системы
\begin{equation}
    \label{eq:2_1:weakOperational}
    A_1 \theta + b \kappa_a (| \theta | \theta^3 - \varphi ) =
    f, A_2 \varphi + \kappa_a (\varphi - |\theta|\theta^3) + F(\varphi, u) = g.
\end{equation}


Для поставленной экстремальной задачи доказано существование решения,
а также выведена система оптимальности, на котором
основан численный алгоритм, представленный в~4.2.2.


Раздел~2.2 содержит постановку задачи без краевых условия для интенсивности излучения.
Будем предполагать, что на границе $\Gamma = \partial \Omega$ известно температурное поле, и тепловой поток:
\begin{equation}
    \label{eq:2_2:bc} \theta = \theta_b, \quad \partial_n\theta = q_b.
\end{equation}

Оптимизационный метода решения краевой задачи~\eqref{eq:2_1:initial},\eqref{eq:2_2:bc}
заключается в рассмотрении задачи граничного
оптимального управления с <<искусственными>> краевыми условиями
\begin{equation}
    \label{eq:2_2:bc3}
    a(\partial_n\theta+\theta) = r,\;\;
    \alpha(\partial_n\varphi+\varphi) = u \text{ на }\Gamma.
\end{equation}
Функция $r(x),\, x\in\Gamma$ является заданной, а неизвестная функция $u(x),\, x\in\Gamma$
играет роль управления.
Экстремальная задача заключается в отыскании тройки
$\{\theta_\lambda,\varphi_\lambda,u_\lambda\}$ такой, что
\begin{equation}
    \label{eq:2_2:cost}
    J_\lambda(\theta, u) = \frac{1}{2}\int\limits_\Gamma (\theta - \theta_b)^2 d\Gamma
    + \frac{\lambda}{2}\int\limits_\Gamma u^2 d\Gamma \rightarrow\inf
\end{equation}
на решениях краевой задачи~\eqref{eq:2_1:initial},\eqref{eq:2_2:bc3}.
Функция $\theta_b(x),\, x\in\Gamma$  и параметр регуляризации $\lambda>0$ заданы.

Для формулировки задачи оптимального управления определим оператор
ограничений $F(\theta, \varphi, u) : V \times V \times U \rightarrow V' \times V'$
\[
    F(\theta, \varphi, u) = \{ aA\theta + b \kappa_a ( [\theta]^4- \varphi) - Br,\;
    \alpha A \varphi + \kappa_a (\varphi -[\theta]^4) - Bu\},
\]
где
\[
    (Ay,z) = (\nabla y, \nabla z) +
    \int\limits_{\Gamma}yz d\Gamma, \;\;\; (Bw, z)
    = \int\limits_{\Gamma}wz d\Gamma.
\]
Тогда задача оптимального управления $(CP)$ заключается в отыскании тройки
$\{\theta, \varphi, u \} \in V \times V \times U$
такой, что
\begin{equation}
    \label{eq:2_2:cp}
    J_\lambda(\theta, u) \equiv \frac{1}{2}\|\theta -\theta_b\|^2_\Gamma
    + \frac{\lambda}{2}\|u\|^2_\Gamma \rightarrow \inf,\;\; F(\theta, \varphi, u)=0.
\end{equation}


Доказана однозначная разрешимость краевой задачи~\eqref{eq:2_1:initial},\eqref{eq:2_2:bc3}
и даны априорные оценки решения, которые далее используются для доказательства разрешимости
задачи $(CP)$.

Доказана следующая теорема
\begin{theorem*}
    Пусть выполняются условия
    $a,b,\alpha,\kappa_a, \lambda ={\textrm Const}> 0$,
    $\theta_b, \,q_b \in U,\;\; r=a(\theta_b+q_b)$
    и существует решение
    задачи~\eqref{eq:2_1:initial},\eqref{eq:2_2:bc}.
    Если $\{\theta_\lambda,\varphi_\lambda,u_\lambda\}$ -- решение
    задачи $(CP)$ для $\lambda>0$, то существует последовательность $\lambda\to +0$
    такая, что
    \[
        \theta_\lambda\rightarrow\theta_*, \;\; \varphi_\lambda\rightarrow\varphi_*
        \text{ слабо в }V,\text{ сильно в }H,
    \]
    где $\theta_*,\varphi_*$ -- решение задачи~\eqref{eq:2_1:initial},\eqref{eq:2_2:bc}.
\end{theorem*}

В разделе~2.3 представлен анализ оптимизационного метода для квазистационарной
модели.
Квазистационарный радиационный и диффузионный теплообмен в ограниченной области
$\Omega \subset \mathbb{R}^{3}$ с границей $\Gamma=\partial \Omega$ моделируем
следующей начально-краевой задачей:

\begin{equation}
    \label{eq:2_3:1}
    \begin{split}
        & \frac{\partial \theta}{\partial t} - a \Delta \theta
        + b \kappa_{a} \left(|\theta| \theta^{3}-\varphi\right) = 0,\\
        & - \alpha \Delta \varphi
        + \kappa_{a} \left(\varphi-|\theta| \theta^{3}\right) = 0,
        \quad x \in \Omega, \quad 0 < t < T;
    \end{split}
\end{equation}
\begin{align}
    a \left(\partial_{n} \theta+\theta\right)=r,
    & \quad \alpha\left(\partial_{n} \varphi
    + \varphi\right) = u \text { на } \Gamma;  \label{eq:2_3:2}\\
    & \left.\theta\right|_{t=0} = \theta_{0}. \label{eq:2_3:3}
\end{align}


Экстремальная задача состоит в том, чтобы найти тройку
$\left\{\theta_{\lambda}, \varphi_{\lambda}, u_{\lambda}\right\}$ такую, что
\begin{equation}
    \label{eq:2_3:4}
    J_{\lambda}(\theta, u)=\frac{1}{2} \int_{0}^{T}
    \int_{\Gamma}\left(\theta-\theta_{b}\right)^{2} d \Gamma d t+\frac{\lambda}{2}
    \int_{0}^{T} \int_{\Gamma} u^{2} d \Gamma d t \rightarrow \inf
\end{equation}
на решениях задачи~\eqref{eq:2_3:1}--\eqref{eq:2_3:3}.

Ставится задача оптимального управления $(OC)$, доказывается существование решение задачи $(OC)$
и выводится система оптимальности.
Представлено доказательство сходимости решений задачи оптимального
управления к решению задачи оптимального управления.


В разделе~2.4 рассмотрена задача сложного теплообмена с условиями Коши для
температуры на части границы.
В данном случае граница области состоит из двух участков,
$\Gamma \coloneqq \partial \Omega =\overline{\Gamma}_1 \cup \overline{\Gamma}_2$,
так что $\Gamma_1 \cap \Gamma_2 =  \emptyset$.
На всей границе $\Gamma$ задается тепловой поток $q_b$,
в качестве условия
переопределения на $\Gamma_1$, в дополнение к условию на
$\varphi$, задается температурное поле $\theta_b$,
\begin{equation}
    \label{eq:2_4:bc2}
    \alpha\partial_n\varphi + \gamma (\varphi - \theta_b ^4 ) = 0,\;
    \theta=\theta_b\quad x\in \Gamma_1.
\end{equation}


Для постановки задачи управления вводится новая неизвестная функция
$\psi= a\theta + \alpha b \varphi$.
Полученная краевая задача имеет вид
\begin{equation}
    \label{eq:2_4:eq2}
    - a \Delta \theta + g (\theta) = \frac{\kappa_a}{\alpha}\psi, \quad
    \Delta \psi = 0, \; x \in \Omega,
\end{equation}
\begin{equation}
    \label{eq:2_4:bc3}
    a \partial_n \theta = q_b, \; \text{ на }\Gamma, \;\;
    \alpha \partial_n \psi + \gamma \psi  =  r,\;\;
    \theta = \theta_b  \text{ на }\Gamma_1.
\end{equation}
Здесь $g(\theta) = b \kappa_a|\theta|\theta^3 + \frac{a\kappa_a}{\alpha}\theta$, $r=\alpha b \gamma \theta_b^4+ \alpha q_b + a \gamma \theta_b$.


Для соответствующей задачи оптимального управления
доказано существование и единственность решения.
Приведено доказательство, что решения задачи оптимального управления
аппроксимируют решение начальной краевой задачи.


\underline{\textbf{Третья глава}} посвящена исследованию квазилинейных моделей.
В разделе~3.1 рассмотрена задача оптимального управления для
квазилинейных уравнений радиационно-кондуктивного
теплообмена, моделирующих процесс внутривенной
лазерной абляции в ограниченной области $\Omega$ с отражающей границей $\Gamma=\partial\Omega$.
$\Omega$ с отражающей границей $\Gamma=\partial\Omega$.
Задача заключается в минимизации функционала
\[ J(\theta)=\int_{0}^{T} \int_{G_{1}}\left(\theta-\theta_{d}\right)^{2} dx dt \rightarrow \inf \]
на решениях начально-краевой задачи:
\begin{equation}
    \label{eq:3_2:1}
    \begin{gathered}
        \sigma \partial \theta / \partial t-\operatorname{div}(k(\theta)
        \nabla \theta)-\beta \varphi=u_{1} \chi \\
        -\operatorname{div}(\alpha \nabla \varphi)+\beta \varphi=u_{2}
        \chi, \quad x \in \Omega, \quad t \in(0, T),
    \end{gathered}
\end{equation}
\begin{equation}
    \label{eq:3_2:2}
    \theta=\left.0\right|_{\Gamma},
    \quad \alpha \partial_{n} \varphi
    +\left.2^{-1} \varphi\right|_{\Gamma}=0,
    \left.\quad \theta\right|_{t=0}=\theta_{0}.
\end{equation}
При этом учитываются ограничения:
\[ u_{1,2} \geq 0, \quad u_{1}+u_{2} \leq P, \left.\quad \theta\right|_{G_{2}} \leq \theta_{*} \]
Здесь $G_{1}$ и $G_{2}$ подмножества $\Omega, \theta$
представляют разницу между реальной температурой
и температурой на границе, которая является постоянной.

Доказано существование решения соответствующей задачи оптимального
управления и даны априорные оценки решений.
Показано, что решения задачи со штрафом сходятся к решению задачи оптимального управления.

В разделе~3.2 приведён анализ метода штрафных функций используемого
для решения задачи оптимального управления с финальным наблюдением.


В \underline{\textbf{четвертой главе}} приведено описание

\FloatBarrier
\pdfbookmark{Заключение}{conclusion}                                  % Закладка pdf
В \underline{\textbf{заключении}} приведены основные результаты работы, которые состоят в следующем:
%% Согласно ГОСТ Р 7.0.11-2011:
%% 5.3.3 В заключении диссертации излагают итоги выполненного исследования, рекомендации, перспективы дальнейшей разработки темы.
%% 9.2.3 В заключении автореферата диссертации излагают итоги данного исследования, рекомендации и перспективы дальнейшей разработки темы.
В диссертации доказано существование квазирешения задачи нахождения
коэффициента отражения участка границы для стационарной модели,
по дополнительной информации о температурном поле.
Экспериментально поверена устойчивость получаемых
решений методом градиентного спуска.
Таким образом, получены важные с теоретической точки зрения результаты,
которые могут быть полезны при дальнейшем использовании
стационарных моделей сложного теплообмена и анализе обратных
задач в рамках нестационарных моделей сложного теплообмена.
Развитые методы исследования начально-краевых задач могут
применяться для изучения различных моделей, описываемых нелинейными
уравнениями со сходной структурой.

Разработанный комплекс программ для постановки численных экспериментов
показал свою надёжность и может в дальнейшем быть использован как
пример для решения подобных задач.

Разработан комплекс программ для проведения вычислительных экспериментов.
Для презентации результатов расчётов, помимо самих солверов,
был разработан программный комплекс для отрисовки
полученных расчётов в трёхмерных областях.

Исследование нестационарных моделей сложного теплообмена и соответствующих им
обратных задач является крайне перспективной областью математического моделирования
и в то же время достаточно сложной для теоретического
анализа и реализации численных решений.


\pdfbookmark{Литература}{bibliography}                                % Закладка pdf
При использовании пакета \verb!biblatex! список публикаций автора по теме
диссертации формируется в разделе <<\publications>>\ файла
\verb!common/characteristic.tex!  при помощи команды \verb!\nocite!

\ifdefmacro{\microtypesetup}{\microtypesetup{protrusion=false}}{} % не рекомендуется применять пакет микротипографики к автоматически генерируемому списку литературы
\urlstyle{rm}                               % ссылки URL обычным шрифтом
\ifnumequal{\value{bibliosel}}{0}{% Встроенная реализация с загрузкой файла через движок bibtex8
    \renewcommand{\bibname}{\large \bibtitleauthor}
    \nocite{*}
    \insertbiblioauthor           % Подключаем Bib-базы
%\insertbiblioexternal   % !!! bibtex не умеет работать с несколькими библиографиями !!!
}{% Реализация пакетом biblatex через движок biber
% Цитирования.
%  * Порядок перечисления определяет порядок в библиографии (только внутри подраздела, если `\insertbiblioauthorgrouped`).
%  * Если не соблюдать порядок "как для \printbibliography", нумерация в `\insertbiblioauthor` будет кривой.
%  * Если цитировать каждый источник отдельной командой --- найти некоторые ошибки будет проще.
%
%% authorvak
    \nocite{vakbib1}%
    \nocite{vakbib2}%
    %
    %% authorwos
    \nocite{wosbib1}%
    %
    %% authorscopus
    \nocite{scbib1}%
    %
    %% authorpathent
    \nocite{patbib1}%
    %
    %% authorprogram
    \nocite{progbib1}%
    %
    %% authorconf
    \nocite{confbib1}%
    \nocite{confbib2}%
    %
    %% authorother
    \nocite{bib1}%
    \nocite{bib2}%

    \ifnumgreater{\value{usefootcite}}{0}{
        \begin{refcontext}[labelprefix={}]
            \ifnum \value{bibgrouped}>0
            \insertbiblioauthorgrouped    % Вывод всех работ автора, сгруппированных по источникам
            \else
            \insertbiblioauthor      % Вывод всех работ автора
            \fi
        \end{refcontext}
    }{
        \ifnum \totvalue{citeexternal}>0
        \begin{refcontext}[labelprefix=A]
            \ifnum \value{bibgrouped}>0
            \insertbiblioauthorgrouped    % Вывод всех работ автора, сгруппированных по источникам
            \else
            \insertbiblioauthor      % Вывод всех работ автора
            \fi
        \end{refcontext}
        \else
        \ifnum \value{bibgrouped}>0
        \insertbiblioauthorgrouped    % Вывод всех работ автора, сгруппированных по источникам
        \else
        \insertbiblioauthor      % Вывод всех работ автора
        \fi
        \fi
    %  \insertbiblioauthorimportant  % Вывод наиболее значимых работ автора (определяется в файле characteristic во второй section)
        \begin{refcontext}[labelprefix={}]
            \insertbiblioexternal            % Вывод списка литературы, на которую ссылались в тексте автореферата
        \end{refcontext}
    % Невидимый библиографический список для подсчёта количества внешних публикаций
    % Используется, чтобы убрать приставку "А" у работ автора, если в автореферате нет
    % цитирований внешних источников.
        \printbibliography[heading=nobibheading, section=0, env=countexternal, keyword=biblioexternal, resetnumbers=true]%
    }
}
\ifdefmacro{\microtypesetup}{\microtypesetup{protrusion=true}}{}
\urlstyle{tt}                               % возвращаем установки шрифта ссылок URL
      % Содержание автореферата

%%% Выходные сведения типографии
\newpage\thispagestyle{empty}

\vspace*{0pt plus1fill}

\small
\begin{center}
    \textit{\thesisAuthor}
    \par\medskip

    \thesisTitle
    \par\medskip

    Автореф.\ дис.\ на соискание ученой степени \thesisDegreeShort
    \par\bigskip

    Подписано в печать \blank[\widthof{999}].\blank[\widthof{999}].\blank[\widthof{99999}].
    Заказ № \blank[\widthof{999999999999}]

    Формат 60\(\times\)90/16.\ Усл.\ печ.\ л.\ 1.\ Тираж 100 экз.
    %Это не совсем формат А5, но наиболее близкий, подробнее: http://ru.wikipedia.org/w/index.php?oldid=78976454

    Типография \blank[0.5\linewidth]
\end{center}
\cleardoublepage

\end{document}
